%% 第22回編集委員会（物理第2巻）：事務局
%% 物理シリーズ　第2巻　熱学　本文（ph1-body.tex と同じ形）
%% 章は \bkinputchap{p2-NN} で読む（熱-0 が p2-00、熱-16 が p2-16。まだない章は、組むときに飛ばして「BKPAGE missing」を出す）
%% 部は \phPartPages{番号}{問い}{デルタのせりふ}{この部を読むと…}：部扉（左）・この部で学ぶこと（右。ph2-part<番号>.tex）・柱とつめ
%%   部扉の3つの問いは、その部の章の「なぜ？」（もくじの問い）から選んだ。裏表紙の5つ（熱-9・3・6・11・7）とは、なるべく重ねない。
% 組版のメモリ対策（全巻にかかる修正事項 No.58。第9巻と同じ）：章の前に Lua のごみ集めをする
\providecommand\bkinputchap[1]{\directlua{collectgarbage("collect") texio.write_nl("log","BKMEM #1 "..math.floor(collectgarbage("count")/1024).." MB")}\IfFileExists{#1.tex}{\input{#1}}{\clearpage\typeout{BKPAGE missing #1}}}
\bkinputchap{p2-00}

%% ======================== 第1部 ふくらむ：熱と第一法則 ========================
\phPartPages{1}%
  {172/{氷水は、氷がとけきるまで0\,℃のままなのはなぜ？},
   160/{自転車の空気入れの筒の先は、なぜ熱くなる？},
   148/{ロケットのノズルは、なぜ一度細くしてから広げる？}}%
  {熱は目に見えないけれど、入った量と出た量は数えられるよ。熱と仕事をまとめて数えれば、空気入れもロケットのノズルも、同じ考え方で計算できるよ。}%
  {温度と熱を区別し、比熱と潜熱で温まり方・とけ方を数え、第一法則で熱と仕事をまとめて数え、理想気体の状態変化と流れる気体のエンタルピーで、ふくらむ袋やノズルの速さを計算できるようになります。この部の考え方は、第2部から先のすべての章の土台です。}
\bkinputchap{p2-01}
\bkinputchap{p2-02}
\bkinputchap{p2-03}
\bkinputchap{p2-04}

%% ======================== 第2部 めぐる：第二法則とサイクル ========================
\phPartPages{2}%
  {172/{エンジンは、燃料の熱を全部、動きに変えられる？},
   160/{高い山の上では、なぜ水が100\,℃より低い温度でわく？},
   148/{圧力なべの中では、水は何℃でわく？}}%
  {熱を全部、仕事に変えることはできないよ。どこまで変えられるかを決めるのが第二法則とエントロピー。エンジンも発電所もエアコンも、熱がめぐるサイクルで考えるよ。}%
  {熱機関の効率の限界をカルノーサイクルで求め、エントロピーで「もとにもどらない」ことを数で表し、エンジン・蒸気の発電所・エアコン・湿り空気を線図と蒸気表で計算し、わく温度と圧力の関係を自由エネルギーで読めるようになります。}
\bkinputchap{p2-05}
\bkinputchap{p2-06}
\bkinputchap{p2-07}
\bkinputchap{p2-08}
\bkinputchap{p2-09}
\bkinputchap{p2-10}

%% ======================== 第3部 移る：伝熱 ========================
\phPartPages{3}%
  {172/{二重窓にすると、なぜ部屋が冷えにくい？},
   160/{パソコンは、なぜファンで風を送って冷やす？},
   148/{たき火に向いた顔だけが熱くなるのはなぜ？}}%
  {熱は、温度の高い所から低い所へ、伝導・対流・ふく射の3つの道で移るよ。壁は熱抵抗、風は熱伝達率、たき火は光。3つを組み合わせると、家もパソコンも設計できるよ。}%
  {熱伝導と熱抵抗で壁や窓から逃げる熱を求め、熱方程式で温度の変わり方を調べ、対流の熱伝達率を相似則と相関式で見積もり、ふく射と熱交換器で熱のやりとりを設計できるようになります。}
\bkinputchap{p2-11}
\bkinputchap{p2-12}
\bkinputchap{p2-13}

%% ======================== 第4部 ぶつかる：ミクロから見た熱 ========================
%% 部の名前が長く、うすい大きな「4」に「熱」の字がかかるので、この巻の第4部だけ、数字を右へ6 mm 寄せる（共通部品は変えない）
\makeatletter\def\bk@partnumpos{158,212}\makeatother
\phPartPages{4}%
  {172/{空気の分子は、どれくらいの速さで飛んでいる？},
   160/{トランプは、切るほどなぜ順番がばらばらになる？},
   148/{電気は、どの方法でつくるとよい？　何をくらべる？}}%
  {温度の正体は、目に見えない分子の動きだよ。分子のぶつかり方と並び方の数を数えると、圧力も比熱もエントロピーも、ミクロから説明できるよ。}%
  {分子の運動から圧力と温度の意味を導き、速さの分布と平均自由行程で粘性や熱伝導を説明し、並び方の数からエントロピーとボルツマン分布を導き、発電の方法をエネルギーの数え方でくらべられるようになります。}
\bkinputchap{p2-14}
\bkinputchap{p2-15}
\bkinputchap{p2-16}

%% ======================== 巻末 ========================
%% 第22回編集委員会（物理第2巻）：事務局
%% 物理シリーズ　第2巻　熱学　巻末（こたえ・記号カード・蒸気表・空気線図・物性値・なぜ？の道すじ・チェック表・つぎの巻へ・さくいん・奥付）
%% 第1巻の ph1-back.tex と同じ部品・同じ形。見出し（\bkbacktitle）の名前は、もくじ（ph2-toc.tex）の巻末の項目の名前と同じにする（tools/check_pages.py p2 が照らし合わせる）。
%% 蒸気表・冷媒の表・湿り空気の表・空気線図・気体と水の物性の表は tools/mksteam2.py が作る（ph2-steamT/P/S・ph2-refrig・ph2-moist・ph2-psychro・ph2-fluid.tex）。手で直さない。
%%   照らし合わせの結果：checks/物理第2巻/巻末/mksteam2_check.txt（本文の状態点 246 点と表の値の食いちがいは 0）。
%% 物性値の出典は、各章のファイルの頭のコメント（出典の欄）と docs/作業メモ/物理第2巻_組どうしの確かめ.md の「巻末の表」の行から集めた。
%%   〔要確認〕は、理科年表・便覧・規格・法令・公式資料の本文そのものでは、まだ確かめていない値（第20回・第22回の決定10）。
\makeatletter
\def\phbkhead#1{\par\addvspace{4mm}\Needspace{6\baselineskip}\noindent{\sffamily\bfseries\large\color{bkDeep}#1}\par\nobreak\vspace{2mm}}
\newcommand\ph@tblstyle{row{1}={bg=bkVol!14, font=\sffamily\bfseries\footnotesize}, hlines={0.4pt,bkLine}, rowsep=1.4pt, colsep=3.5pt}
\newcommand\phyk{〔要確認〕}
%% ---------------- こたえ ----------------
\bkAnswerPages
\par\vspace{6mm}
\begin{tcolorbox}[bk base, colframe=bkTheme!60, colback=bkTheme!5, boxrule=0.8pt, arc=2mm, left=4mm, right=4mm, top=2.5mm, bottom=3mm]
{\sffamily\bfseries\color{bkDeep}答え合わせのしかた}\par\vspace{1.5mm}
\begin{enumerate}[label={\bkcircled{\arabic*}}, leftmargin=2.4\zw, itemsep=1pt, topsep=0pt]
\item まず自分の答えを全部書いてから、ここを見よう。
\item 数だけでなく、単位と、桁（10倍・100倍のちがい）を確かめよう。温度は、式の中ではいつも K です。℃ のまま比やかけ算をしていないかも確かめよう。
\item 蒸気表・空気線図を読む問題は、表の読み方（補間のしかた）で、最後の桁が1つずれることがあります。それは間違いではありません。
\item ちがっていたら、同じ形の例題にもどって、どの段（何を無視するか・式を立てる・解く）でちがったかを探そう。熱と仕事の向き（入る・出る）も確かめよう。
\end{enumerate}
\end{tcolorbox}

%% ---------------- 記号カード ----------------
%% 数学シリーズの記号で、この巻でよく使うもの（この巻の章の \bkkigocard は自動でうしろに集まる）
\bkKigoPages{%
\bkkigorow{$\Delta$}{デルタ}{変化量。「あと$-$まえ」。温度差 $\Delta T$〔K〕、内部エネルギーの変化 $\Delta U$}{2-3}%
%% 修2-6（U-6・Q6-10）：$\dot Q$ は熱-0 のカード（キュー・ドット）だけに残し、この行からけずった
\bkkigorow{$\tfrac{\dd T}{\dd t}$}{ディーティー・ディーティー}{時間で微分。温度の変わる速さ（熱-1 の冷め方、熱-11）}{5-5}%
\bkkigorow{$\approx$}{ニアリーイコール}{ほぼ等しい（近似）}{5-8}%
\bkkigorow{$\displaystyle\int$}{インテグラル}{細かく分けてたす（積分）。仕事 $\int p\,\dd V$}{5-10}%
\bkkigorow{$\displaystyle\sum$}{シグマ}{全部たす。分配関数 $Z=\sum e^{-\varepsilon_i/k_{\mathrm B}T}$（熱-15）}{5-10}%
\bkkigorow{$e$、\ $\exp$}{イー、エクスポネンシャル}{ネイピア数。$\exp x=e^x$。ボルツマン因子、冷め方}{5-14}%
\bkkigorow{$\ln x$}{エル・エヌ・エックス}{自然対数。等温の仕事 $nRT\ln\frac{V_2}{V_1}$、エントロピー}{5-14}%
\bkkigorow{$n!$}{エヌのかいじょう}{1から $n$ までの整数をすべてかけた数（熱-15 の並び方の数）}{5-17}%
\bkkigorow{$\langle v^2\rangle$、\ $\bar v$}{へいきん}{平均。分子の速さの平均（熱-14）}{3-13}%
\bkkigorow{$\tfrac{\partial f}{\partial x}$}{ラウンドディー}{偏微分。ほかの変数を止めて $x$ だけ動かす}{8-2}%
\bkkigorow{$\bigl(\tfrac{\partial U}{\partial T}\bigr)_V$}{かっこの添字}{止めておく変数を添字で書く。定積比熱（熱-2・熱-10）}{8-3}%
\bkkigorow{$\dd U$、\ $\delta Q$}{ディー・ユー、デルタ・キュー}{$\dd U$ は全微分（状態で決まる）。$\delta Q$・$\delta W$ は道で変わる小さな量}{8-3}%
\bkkigorow{$\oint$}{しゅうかいせきぶん}{閉じた道を1周たす。サイクルの仕事・$\oint\dd U=0$}{8-7}%
\bkkigorow{$\mathrm{grad}$、\ $\nabla$}{グラディエント、ナブラ}{いちばん急な向き。熱流束 $\vec q=-k\,\mathrm{grad}\,T$}{8-4}%
\bkkigorow{$\mathrm{div}$}{ダイバージェンス}{わき出し。熱の流れの袋（熱-11）}{8-9}%
\bkkigorow{$\nabla^2$}{ラプラシアン}{ならす働き。熱方程式 $\tfrac{\partial T}{\partial t}=\alpha\nabla^2T$}{8-12}%
}

%% 第1巻はここで改ページしたが、記号カードの最後のページに大きな空きが残るので、続けて組む（見出しは \Needspace で次のページへ送る）
\phbkhead{分野で記号がぶつかるもの}
{\small 熱学は、使う量がとても多く、同じ文字がちがう量を表すことがあります。この本では、同じ章の中では1つの文字を1つの意味にだけ使い、ぶつかるときは文字や添字を変えました。どの量かは、章と単位で見分けよう。右から2つめの列は、第1巻『力学』とほかの巻での意味です。\par}\vspace{2mm}
{\footnotesize
\begin{longtblr}{colspec={Q[c,m,wd=15mm] X[l,m] Q[l,m,wd=34mm] Q[l,m,wd=28mm]}, width=\linewidth, rowhead=1, \ph@tblstyle}
記号 & この巻での意味 & 第1巻・ほかの巻での意味 & 見分け方 \\
$T$、$\theta$ & 絶対温度 $T$〔K〕（全章。式の中の $T$ はいつも K）。セルシウス温度 $\theta$〔℃〕（熱-1）、アインシュタイン温度 $\theta_{\mathrm E}$（熱-15）。周期は熱学では文字にしない & 周期・ねじりモーメント $T$（力-1・力-11）、角 $\theta$（4-2）、数学第6巻の温度 $T$〔℃〕 & 単位：K、℃ \\
$Q$、$\dot Q$、$q$ & 熱量 $Q$（熱-1）、熱の流れ $\dot Q$（熱-11）、熱流束 $q$〔W/m$^2$〕（熱-11）、1\,kg あたりの熱 $q$（熱-4・熱-7・熱-8） & 体積流量・せん断力 $Q$（力-13・力-10）、電荷（電-1） & 単位：J、W、W/m$^2$、J/kg \\
$W$、$w$ & される仕事 $W$・する仕事 $W'$（熱-2）、サイクルの $W_{\text{出}}$・$W_{\text{入}}$（熱-5）、1\,kg あたり $w$（熱-4・熱-7）。重さは $mg$ & 仕事（力-4） & 単位：J、J/kg \\
$\Omega$ & 状態の数（場合の数、熱-15。専門書の $S=k\log W$ の $W$）、立体角 $\dd\Omega$（熱-13） & 電気抵抗の単位 Ω（電-4） & 状態の数は単位なし、立体角は sr。単位の Ω は立体 \\
$U$ & 内部エネルギー（熱-2）。熱貫流率・熱通過率 $U$（熱-11・熱-13 だけ） & ― & 単位：J、W/(m$^2$·K) \\
$h$、$H$ & 比エンタルピー $h$（熱-4・熱-8・熱-9）、熱伝達率 $h$（熱-11〜熱-13）、エンタルピー $H$（熱-4）、大気の高さの目安 $H$（熱-3・熱-15）。高さは $z$ & 高さ・落差 $h$、全ヘッド $H$（力-4・力-13）、プランク定数（原-2） & 単位：J/kg、W/(m$^2$·K)、J、m \\
$s$、$S$ & 比エントロピー $s$・エントロピー $S$（熱-6・熱-8） & 動いた道のり（力-1・力-4） & 単位：J/(kg·K)、J/K \\
$k$ & 熱伝導率（熱-11〜熱-14）、ポリトロープ指数（熱-3）。ボルツマン定数はいつも $k_{\mathrm B}$、ばね定数は $k_{\mathrm s}$（熱-14・熱-15） & ばね定数（力-2）、波数（波-5） & 単位：W/(m·K)、単位なし \\
$c$、$C$ & 比熱 $c$・$c_p$・$c_v$（熱-1・熱-3）、モル比熱 $C_V$・$C_p$（熱-3）、熱容量 $C_{\mathrm{th}}$（熱-1・熱-11）、熱容量流量 $C$（熱-13）。光の速さは $c_0$ & ダンパーの係数 $c$（力-2）、静電容量 $C$（電-3） & 単位：J/(kg·K)、J/(mol·K)、J/K、W/K \\
$R$ & モル気体定数 $R$（熱-3）、比気体定数 $R_{\mathrm s}=R/M$（熱-3）、熱抵抗 $R_{\mathrm{th}}$（熱-11）。機械の本は1\,kg あたりを $R$ と書く & 垂直抗力（力-5）、電気抵抗（電-4） & 単位：J/(mol·K)、J/(kg·K)、K/W \\
$x$ & 乾き度（熱-8）、絶対湿度（熱-9） & 位置（力-1） & 単位なし、kg/kg(DA) \\
$v$ & 比体積（熱-8・熱-10）、流れの速さ（熱-4）、分子の速さ（熱-14）。熱-4 では比体積を $1/\rho$ と書く & 速度（力-1） & 単位：m$^3$/kg、m/s \\
$\varepsilon$ & 圧縮比（熱-7）、放射率（熱-13）、エネルギーの準位 $\varepsilon_i$（熱-15） & ひずみ（力-8）、誘電率（電-3） & 圧縮比・放射率は単位なし \\
$\sigma$ & ステファン・ボルツマン定数（熱-13）、締切比（熱-7）、表面張力（熱-12） & 応力（力-8）、標準偏差（3-14）、導電率（電-4） & 単位：W/(m$^2$·K$^4$)、単位なし、N/m \\
$\alpha$ & 線膨張係数（熱-1）、熱拡散率（熱-11・熱-12）、吸収率（熱-13）、未定乗数（熱-15） & 線膨張係数（力-8）、次元解析の指数（力-15） & 単位：1/K、m$^2$/s、単位なし \\
$\beta$ & 体膨張率（熱-1・熱-12）、$\beta=1/(k_{\mathrm B}T)$（熱-15） & ― & 単位：1/K、1/J \\
$\gamma$ & 比熱比 $C_p/C_V$（熱-3）。機械の本では $\kappa$ & 次元解析の指数（力-15） & 単位なし \\
$\mu$ & 粘度（熱-12・熱-14）、化学ポテンシャル（熱-10・熱-15） & 摩擦係数・粘度（力-2・力-14）、透磁率（電-7） & 単位：Pa·s、J/mol \\
$\lambda$ & 空気比（熱-1）、波長（熱-13）、平均自由行程（熱-14）。JIS・建築の本では熱伝導率も $\lambda$ & 管摩擦係数（力-14）、固有値（7-8） & 空気比は単位なし、m \\
$L$、$l$ & 潜熱 $L$（熱-1・熱-8・熱-10・熱-12）、壁の厚さ $L$（熱-11）、代表の長さ $l$（熱-12） & 長さ（力-8）、インダクタンス（電-8） & 単位：J/kg、m \\
$n$ & 物質量（熱-3）、回転数（熱-7）、数密度 $n_V$（熱-14）、棚にいる数 $n_i$（熱-15） & ― & 単位：mol、min$^{-1}$、1/m$^3$ \\
$Z$ & 分配関数（熱-15）、圧縮係数（熱-3・熱-10） & インピーダンス（電-10） & どちらも単位なし。章で見分ける \\
$F$、$G$ & ヘルムホルツ・ギブズの自由エネルギー（熱-10）、形態係数 $F_{12}$（熱-13） & 力 $F$（力-2）、万有引力定数・横弾性係数 $G$（力-2・力-8） & 単位：J、単位なし \\
$P$、$p$ & 圧力 $p$（いつも絶対圧）、仕事率・電力 $P$（熱-1・熱-2）、確率 $P_i$（熱-15）、フィンの周りの長さ $P$（熱-12） & 運動量 $\vec p$（力-3） & 単位：Pa、W \\
$a$ & 音速（熱-4）、ファンデルワールスの定数 $a$・$b$（熱-10） & 加速度（力-1）、水撃の波の速さ（力-14） & 単位：m/s \\
$\tau$、$\rho$ & 時定数 $\tau$（熱-11・熱-12）、透過率 $\tau$・反射率 $\rho$（熱-13）、密度 $\rho$ & せん断応力 $\tau$（力-8）、抵抗率 $\rho$（電-4） & 単位：s、単位なし、kg/m$^3$ \\
$\delta$ & 道で変わる小さな量 $\delta Q$（熱-2）、境界層の厚さ $\delta$（熱-12） & ― & 単位：J、m \\
$\eta$ & 効率（熱-5・熱-7・熱-8）、熱交換器の温度効率 $\eta_T$（熱-13）。工学の本の温度効率 $\varepsilon$ は使わない & ― & 単位なし \\
$\varphi$ & 相対湿度（熱-9） & 球座標の角（4-15）、電位 $\phi$（8-4） & 単位：\% \\
$m$ & 質量。フィンの式の $m=\sqrt{hP/(kA_{\mathrm c})}$（熱-12） & 質量（力-2） & 単位：kg、1/m \\
\end{longtblr}\par}

\phbkhead{この巻の約束（記号と向き）}
{\small
\begin{itemize}[leftmargin=1.2\zw, itemsep=2pt, topsep=1pt, label={\color{bkVol}●}]
\item \textbf{第一法則}：本文は $\Delta U=Q+W$（$Q$ は受けとる熱、$W$ は\textbf{される}仕事。高校の教科書と同じ）。機械の本の $Q=\Delta U+W$（$W$ は\textbf{する}仕事）とは $W$ の向きが反対です（熱-2）。
\item \textbf{サイクルの熱と仕事}：大きさ（正の数）と向きの添字で書きます。$Q_{\mathrm H}$・$Q_{\mathrm L}$、$W_{\text{出}}$・$W_{\text{入}}$、1\,kg あたりは $q_{\text{入}}$・$q_{\text{出}}$・$w_{\text{出}}$・$w_{\text{入}}$（熱-5・熱-7・熱-8）。
\item \textbf{成績係数}：本文は COP$_{\text{冷}}$（冷房・冷凍）・COP$_{\text{暖}}$（暖房）。専門書では COP$_{\mathrm R}$・COP$_{\mathrm H}$ と書きます（熱-5）。
\item \textbf{蒸気表の印}：$'$ は飽和液、$''$ は飽和蒸気の値（$h'$・$h''$。英語の本では $h_f$・$h_g$）。蒸発熱は $L=h''-h'$（熱-8）。
\item \textbf{湿り空気}：比エンタルピーと絶対湿度は「乾いた空気\qty{1}{kg}あたり」で、単位に (DA) をつけます（熱-9）。
\item \textbf{圧力}：気体の式には、いつも絶対圧を入れます。ゲージ圧は、ことばで「ゲージ圧」と書きそえます（熱-3）。
\end{itemize}\par}

%% ---------------- 蒸気表・空気線図・物性値 ----------------
\bkbacktitle{蒸気表・空気線図・物性値}
{\small この巻で使った値です。水と水蒸気の表は IAPWS-IF97（国際水・蒸気性質協会の工業用の式）から、冷媒と湿り空気と気体の表は CoolProp 8.0（各物質の状態式）から、Python のプログラム（\texttt{tools/mksteam2.py}）で計算しました。本文の例題の値は、この表と同じ丸めです。そのほかの物性値は温度や材料の種類で変わるので、設計では便覧・規格の値を使います。\phyk は、便覧・規格・法令・公式資料の本文そのものでは、まだ確かめていない値です。\par}\vspace{1mm}
{\footnotesize\color{bkGray}丸め：温度は小数2けた、$h$ は小数1けた、$s$ は小数4けた、$v'$ は小数6けた、$v''$ と過熱蒸気の $v$ は4けた（熱-8 と同じ）。$h$・$s$ の0の点：水は三重点の飽和液（$u'=0$、$s'=0$）、冷媒は\qty{0}{\celsius}の飽和液で $h'=\qty{200}{kJ/kg}$・$s'=\qty{1}{kJ/(kg.K)}$（国際冷凍協会 IIR の約束）。だから、ちがう表の $h$ どうしを引き算してはいけません。\par}

\phbkhead{1　水の飽和表（温度基準）}
{\footnotesize\color{bkGray}\qty{20}{\celsius}の飽和蒸気圧 \qty{2.339}{kPa}（熱-8・熱-9・熱-10、第1巻 力-14）。\textsuperscript{*}は臨界点（\qty{373.95}{\celsius}・\qty{22.064}{MPa}）。臨界点では液と蒸気の区別がなくなり、蒸発熱は0になる（熱-8・熱-10）。\par}\vspace{0.5mm}
%% 自動生成（tools/mksteam2.py）。直接編集しないこと。値は IAPWS-IF97 の式（iapws 1.5.5）・CoolProp 8.0.0
\begin{longtblr}{colspec={Q[c,m,wd=11.5mm] Q[r,m,wd=13mm] Q[r,m] Q[r,m] Q[r,m] Q[r,m] Q[r,m] Q[r,m] Q[r,m]}, width=\linewidth, rowhead=1, colsep=3.2pt, rowsep=0.35pt, hlines={0.3pt,bkLine}, vline{2}={0.4pt,bkLine}, row{1-1}={bg=bkVol!14, font=\sffamily\bfseries\scriptsize, halign=c}, cells={font=\scriptsize}}
{温度\\$\theta$〔℃〕} & {圧力\\$p_{\mathrm s}$〔kPa〕} & {$v'$\\〔m$^3$/kg〕} & {$v''$\\〔m$^3$/kg〕} & {$h'$\\〔kJ/kg〕} & {$h''$\\〔kJ/kg〕} & {蒸発熱\\$L=h''-h'$} & {$s'$\\〔kJ/(kg·K)〕} & {$s''$\\〔kJ/(kg·K)〕} \\
0.01 & 0.6117 & 0.001000 & 206.0 & 0.0 & 2500.9 & 2500.9 & 0.0000 & 9.1555 \\
2 & 0.7060 & 0.001000 & 179.8 & 8.4 & 2504.6 & 2496.2 & 0.0306 & 9.1027 \\
4 & 0.8135 & 0.001000 & 157.1 & 16.8 & 2508.2 & 2491.4 & 0.0611 & 9.0506 \\
5 & 0.8726 & 0.001000 & 147.0 & 21.0 & 2510.1 & 2489.1 & 0.0763 & 9.0249 \\
6 & 0.9354 & 0.001000 & 137.6 & 25.2 & 2511.9 & 2486.7 & 0.0913 & 8.9994 \\
8 & 1.0730 & 0.001000 & 120.8 & 33.6 & 2515.6 & 2481.9 & 0.1213 & 8.9492 \\
10 & 1.2282 & 0.001000 & 106.3 & 42.0 & 2519.2 & 2477.2 & 0.1511 & 8.8998 \\
12 & 1.4028 & 0.001001 & 93.72 & 50.4 & 2522.9 & 2472.5 & 0.1806 & 8.8514 \\
14 & 1.5989 & 0.001001 & 82.80 & 58.8 & 2526.5 & 2467.7 & 0.2099 & 8.8038 \\
15 & 1.7057 & 0.001001 & 77.88 & 63.0 & 2528.4 & 2465.4 & 0.2245 & 8.7804 \\
16 & 1.8188 & 0.001001 & 73.29 & 67.2 & 2530.2 & 2463.0 & 0.2390 & 8.7571 \\
18 & 2.0647 & 0.001001 & 65.00 & 75.5 & 2533.8 & 2458.3 & 0.2678 & 8.7112 \\
20 & 2.3392 & 0.001002 & 57.76 & 83.9 & 2537.5 & 2453.5 & 0.2965 & 8.6661 \\
22 & 2.6452 & 0.001002 & 51.42 & 92.3 & 2541.1 & 2448.8 & 0.3250 & 8.6218 \\
24 & 2.9856 & 0.001003 & 45.86 & 100.7 & 2544.7 & 2444.1 & 0.3532 & 8.5783 \\
25 & 3.1697 & 0.001003 & 43.34 & 104.8 & 2546.5 & 2441.7 & 0.3673 & 8.5568 \\
26 & 3.3637 & 0.001003 & 40.98 & 109.0 & 2548.4 & 2439.3 & 0.3813 & 8.5355 \\
28 & 3.7828 & 0.001004 & 36.68 & 117.4 & 2552.0 & 2434.6 & 0.4091 & 8.4934 \\
30 & 4.2467 & 0.001004 & 32.88 & 125.7 & 2555.6 & 2429.8 & 0.4368 & 8.4521 \\
32 & 4.7592 & 0.001005 & 29.53 & 134.1 & 2559.2 & 2425.1 & 0.4643 & 8.4115 \\
34 & 5.3247 & 0.001006 & 26.56 & 142.5 & 2562.8 & 2420.3 & 0.4916 & 8.3715 \\
35 & 5.6286 & 0.001006 & 25.21 & 146.6 & 2564.6 & 2417.9 & 0.5052 & 8.3518 \\
36 & 5.9475 & 0.001006 & 23.93 & 150.8 & 2566.4 & 2415.6 & 0.5187 & 8.3323 \\
38 & 6.6324 & 0.001007 & 21.60 & 159.2 & 2570.0 & 2410.8 & 0.5457 & 8.2936 \\
40 & 7.3844 & 0.001008 & 19.52 & 167.5 & 2573.5 & 2406.0 & 0.5724 & 8.2557 \\
45 & 9.5944 & 0.001010 & 15.25 & 188.4 & 2582.5 & 2394.0 & 0.6386 & 8.1634 \\
50 & 12.351 & 0.001012 & 12.03 & 209.3 & 2591.3 & 2382.0 & 0.7038 & 8.0749 \\
55 & 15.761 & 0.001015 & 9.565 & 230.2 & 2600.1 & 2369.9 & 0.7680 & 7.9899 \\
60 & 19.946 & 0.001017 & 7.668 & 251.2 & 2608.8 & 2357.7 & 0.8312 & 7.9082 \\
65 & 25.041 & 0.001020 & 6.194 & 272.1 & 2617.5 & 2345.4 & 0.8935 & 7.8296 \\
70 & 31.201 & 0.001023 & 5.040 & 293.0 & 2626.1 & 2333.1 & 0.9550 & 7.7540 \\
75 & 38.595 & 0.001026 & 4.129 & 314.0 & 2634.6 & 2320.6 & 1.0156 & 7.6812 \\
80 & 47.415 & 0.001029 & 3.405 & 334.9 & 2643.0 & 2308.1 & 1.0754 & 7.6110 \\
85 & 57.867 & 0.001032 & 2.826 & 355.9 & 2651.3 & 2295.4 & 1.1344 & 7.5434 \\
90 & 70.182 & 0.001036 & 2.359 & 377.0 & 2659.5 & 2282.6 & 1.1927 & 7.4781 \\
95 & 84.609 & 0.001040 & 1.981 & 398.0 & 2667.6 & 2269.6 & 1.2502 & 7.4150 \\
100 & 101.42 & 0.001043 & 1.672 & 419.1 & 2675.6 & 2256.5 & 1.3070 & 7.3541 \\
105 & 120.90 & 0.001047 & 1.418 & 440.2 & 2683.4 & 2243.2 & 1.3632 & 7.2951 \\
110 & 143.38 & 0.001052 & 1.209 & 461.4 & 2691.1 & 2229.7 & 1.4187 & 7.2380 \\
115 & 169.18 & 0.001056 & 1.036 & 482.6 & 2698.6 & 2216.0 & 1.4735 & 7.1827 \\
120 & 198.67 & 0.001060 & 0.8913 & 503.8 & 2705.9 & 2202.1 & 1.5278 & 7.1291 \\
125 & 232.22 & 0.001065 & 0.7701 & 525.1 & 2713.1 & 2188.0 & 1.5815 & 7.0770 \\
130 & 270.26 & 0.001070 & 0.6681 & 546.4 & 2720.1 & 2173.7 & 1.6346 & 7.0264 \\
140 & 361.50 & 0.001080 & 0.5085 & 589.2 & 2733.4 & 2144.2 & 1.7393 & 6.9293 \\
150 & 476.10 & 0.001091 & 0.3925 & 632.3 & 2745.9 & 2113.7 & 1.8420 & 6.8370 \\
160 & 618.14 & 0.001102 & 0.3068 & 675.6 & 2757.4 & 2081.9 & 1.9428 & 6.7491 \\
170 & 792.05 & 0.001114 & 0.2426 & 719.2 & 2767.9 & 2048.7 & 2.0419 & 6.6649 \\
180 & 1003 & 0.001127 & 0.1939 & 763.2 & 2777.2 & 2014.0 & 2.1395 & 6.5841 \\
190 & 1255 & 0.001141 & 0.1564 & 807.6 & 2785.3 & 1977.7 & 2.2358 & 6.5060 \\
200 & 1555 & 0.001157 & 0.1272 & 852.4 & 2792.1 & 1939.7 & 2.3308 & 6.4303 \\
210 & 1907 & 0.001173 & 0.1043 & 897.7 & 2797.4 & 1899.6 & 2.4248 & 6.3565 \\
220 & 2319 & 0.001190 & 0.08610 & 943.6 & 2801.1 & 1857.4 & 2.5178 & 6.2842 \\
230 & 2797 & 0.001209 & 0.07151 & 990.2 & 2803.0 & 1812.8 & 2.6102 & 6.2131 \\
240 & 3347 & 0.001229 & 0.05971 & 1037.5 & 2803.1 & 1765.5 & 2.7019 & 6.1425 \\
250 & 3976 & 0.001252 & 0.05009 & 1085.7 & 2801.0 & 1715.3 & 2.7934 & 6.0722 \\
260 & 4692 & 0.001276 & 0.04218 & 1134.8 & 2796.6 & 1661.8 & 2.8847 & 6.0017 \\
270 & 5503 & 0.001303 & 0.03562 & 1185.1 & 2789.7 & 1604.6 & 2.9762 & 5.9304 \\
280 & 6416 & 0.001333 & 0.03015 & 1236.7 & 2779.8 & 1543.2 & 3.0681 & 5.8578 \\
290 & 7442 & 0.001366 & 0.02556 & 1289.8 & 2766.6 & 1476.8 & 3.1608 & 5.7832 \\
300 & 8588 & 0.001404 & 0.02166 & 1344.8 & 2749.6 & 1404.8 & 3.2547 & 5.7058 \\
310 & 9865 & 0.001448 & 0.01834 & 1402.0 & 2727.9 & 1325.9 & 3.3506 & 5.6243 \\
320 & 11\,284 & 0.001499 & 0.01548 & 1462.1 & 2700.7 & 1238.6 & 3.4491 & 5.5373 \\
330 & 12\,858 & 0.001561 & 0.01298 & 1525.7 & 2666.2 & 1140.5 & 3.5516 & 5.4425 \\
340 & 14\,600 & 0.001638 & 0.01078 & 1594.4 & 2622.1 & 1027.6 & 3.6599 & 5.3359 \\
350 & 16\,529 & 0.001740 & 0.008801 & 1670.9 & 2563.6 & 892.7 & 3.7783 & 5.2109 \\
360 & 18\,666 & 0.001895 & 0.006945 & 1761.5 & 2481.0 & 719.5 & 3.9164 & 5.0527 \\
370 & 21\,043 & 0.002222 & 0.004946 & 1892.6 & 2333.5 & 440.9 & 4.1142 & 4.7996 \\
373.95\textsuperscript{*} & 22\,064 & 0.003106 & 0.003106 & 2087.5 & 2087.5 & 0.0 & 4.4120 & 4.4120 \\
\end{longtblr}


\phbkhead{2　水の飽和表（圧力基準）}
{\footnotesize\color{bkGray}0.000612\,MPa は三重点の圧力（\qty{611.657}{Pa}）、0.101325\,MPa は1気圧。0.034\,MPa（エベレスト山頂の気圧の目安）、0.063\,MPa（富士山頂の気圧の目安）は、熱-8・熱-10 の例で使った圧力。\par}\vspace{0.5mm}
%% 自動生成（tools/mksteam2.py）。直接編集しないこと。値は IAPWS-IF97 の式（iapws 1.5.5）・CoolProp 8.0.0
\begin{longtblr}{colspec={Q[c,m,wd=13mm] Q[r,m,wd=12mm] Q[r,m] Q[r,m] Q[r,m] Q[r,m] Q[r,m] Q[r,m] Q[r,m]}, width=\linewidth, rowhead=1, colsep=3.2pt, rowsep=0.35pt, hlines={0.3pt,bkLine}, vline{2}={0.4pt,bkLine}, row{1-1}={bg=bkVol!14, font=\sffamily\bfseries\scriptsize, halign=c}, cells={font=\scriptsize}}
{圧力\\$p$〔MPa〕} & {飽和温度\\$\theta_{\mathrm s}$〔℃〕} & {$v'$\\〔m$^3$/kg〕} & {$v''$\\〔m$^3$/kg〕} & {$h'$\\〔kJ/kg〕} & {$h''$\\〔kJ/kg〕} & {蒸発熱\\$L=h''-h'$} & {$s'$\\〔kJ/(kg·K)〕} & {$s''$\\〔kJ/(kg·K)〕} \\
0.000612 & 0.01 & 0.001000 & 206.0 & 0.0 & 2500.9 & 2500.9 & 0.0000 & 9.1555 \\
0.001 & 6.97 & 0.001000 & 129.2 & 29.3 & 2513.7 & 2484.4 & 0.1059 & 8.9749 \\
0.0015 & 13.02 & 0.001001 & 87.96 & 54.7 & 2524.7 & 2470.1 & 0.1956 & 8.8270 \\
0.002 & 17.50 & 0.001001 & 66.99 & 73.4 & 2532.9 & 2459.5 & 0.2606 & 8.7227 \\
0.0025 & 21.08 & 0.001002 & 54.24 & 88.4 & 2539.4 & 2451.0 & 0.3119 & 8.6422 \\
0.003 & 24.08 & 0.001003 & 45.66 & 101.0 & 2544.9 & 2443.9 & 0.3543 & 8.5766 \\
0.004 & 28.96 & 0.001004 & 34.79 & 121.4 & 2553.7 & 2432.3 & 0.4224 & 8.4735 \\
0.005 & 32.88 & 0.001005 & 28.19 & 137.8 & 2560.8 & 2423.0 & 0.4763 & 8.3939 \\
0.006 & 36.16 & 0.001006 & 23.73 & 151.5 & 2566.7 & 2415.2 & 0.5209 & 8.3291 \\
0.0075 & 40.29 & 0.001008 & 19.23 & 168.8 & 2574.1 & 2405.3 & 0.5763 & 8.2502 \\
0.008 & 41.51 & 0.001008 & 18.10 & 173.9 & 2576.2 & 2402.4 & 0.5925 & 8.2274 \\
0.01 & 45.81 & 0.001010 & 14.67 & 191.8 & 2583.9 & 2392.1 & 0.6492 & 8.1489 \\
0.015 & 53.97 & 0.001014 & 10.02 & 225.9 & 2598.3 & 2372.4 & 0.7548 & 8.0071 \\
0.02 & 60.06 & 0.001017 & 7.648 & 251.4 & 2608.9 & 2357.5 & 0.8320 & 7.9072 \\
0.025 & 64.96 & 0.001020 & 6.203 & 271.9 & 2617.4 & 2345.5 & 0.8931 & 7.8302 \\
0.03 & 69.10 & 0.001022 & 5.229 & 289.2 & 2624.6 & 2335.3 & 0.9439 & 7.7675 \\
0.034 & 72.00 & 0.001024 & 4.650 & 301.4 & 2629.5 & 2328.1 & 0.9793 & 7.7245 \\
0.04 & 75.86 & 0.001026 & 3.993 & 317.6 & 2636.1 & 2318.5 & 1.0259 & 7.6690 \\
0.05 & 81.32 & 0.001030 & 3.240 & 340.5 & 2645.2 & 2304.7 & 1.0910 & 7.5930 \\
0.063 & 87.18 & 0.001034 & 2.610 & 365.1 & 2654.9 & 2289.8 & 1.1599 & 7.5146 \\
0.07 & 89.93 & 0.001036 & 2.365 & 376.7 & 2659.4 & 2282.7 & 1.1919 & 7.4790 \\
0.075 & 91.76 & 0.001037 & 2.217 & 384.4 & 2662.4 & 2278.0 & 1.2130 & 7.4557 \\
0.08 & 93.49 & 0.001038 & 2.087 & 391.6 & 2665.2 & 2273.5 & 1.2328 & 7.4339 \\
0.09 & 96.69 & 0.001041 & 1.869 & 405.1 & 2670.3 & 2265.2 & 1.2694 & 7.3942 \\
0.1 & 99.61 & 0.001043 & 1.694 & 417.4 & 2674.9 & 2257.5 & 1.3026 & 7.3588 \\
0.101325 & 99.97 & 0.001043 & 1.673 & 419.0 & 2675.5 & 2256.5 & 1.3067 & 7.3544 \\
0.125 & 105.97 & 0.001048 & 1.375 & 444.3 & 2684.9 & 2240.6 & 1.3739 & 7.2840 \\
0.15 & 111.35 & 0.001053 & 1.159 & 467.1 & 2693.1 & 2226.0 & 1.4335 & 7.2229 \\
0.175 & 116.04 & 0.001057 & 1.004 & 487.0 & 2700.1 & 2213.2 & 1.4849 & 7.1714 \\
0.2 & 120.21 & 0.001061 & 0.8857 & 504.7 & 2706.2 & 2201.6 & 1.5301 & 7.1269 \\
0.25 & 127.41 & 0.001067 & 0.7187 & 535.4 & 2716.5 & 2181.2 & 1.6072 & 7.0524 \\
0.3 & 133.53 & 0.001073 & 0.6058 & 561.5 & 2724.9 & 2163.4 & 1.6718 & 6.9916 \\
0.35 & 138.86 & 0.001079 & 0.5242 & 584.3 & 2732.0 & 2147.7 & 1.7275 & 6.9401 \\
0.4 & 143.61 & 0.001084 & 0.4624 & 604.7 & 2738.1 & 2133.3 & 1.7766 & 6.8954 \\
0.45 & 147.91 & 0.001088 & 0.4139 & 623.2 & 2743.4 & 2120.2 & 1.8206 & 6.8560 \\
0.5 & 151.84 & 0.001093 & 0.3748 & 640.2 & 2748.1 & 2107.9 & 1.8606 & 6.8206 \\
0.6 & 158.83 & 0.001101 & 0.3156 & 670.5 & 2756.1 & 2085.6 & 1.9311 & 6.7592 \\
0.7 & 164.95 & 0.001108 & 0.2728 & 697.1 & 2762.7 & 2065.6 & 1.9921 & 6.7070 \\
0.8 & 170.41 & 0.001115 & 0.2403 & 721.0 & 2768.3 & 2047.3 & 2.0460 & 6.6615 \\
0.9 & 175.36 & 0.001121 & 0.2149 & 742.7 & 2773.0 & 2030.3 & 2.0944 & 6.6212 \\
1 & 179.89 & 0.001127 & 0.1943 & 762.7 & 2777.1 & 2014.4 & 2.1384 & 6.5850 \\
1.2 & 187.96 & 0.001139 & 0.1632 & 798.5 & 2783.8 & 1985.3 & 2.2163 & 6.5217 \\
1.4 & 195.05 & 0.001149 & 0.1408 & 830.1 & 2788.9 & 1958.8 & 2.2839 & 6.4675 \\
1.5 & 198.30 & 0.001154 & 0.1317 & 844.7 & 2791.0 & 1946.3 & 2.3147 & 6.4431 \\
1.6 & 201.38 & 0.001159 & 0.1237 & 858.6 & 2792.9 & 1934.3 & 2.3438 & 6.4200 \\
1.8 & 207.12 & 0.001168 & 0.1104 & 884.6 & 2796.0 & 1911.4 & 2.3978 & 6.3776 \\
2 & 212.38 & 0.001177 & 0.09958 & 908.6 & 2798.4 & 1889.8 & 2.4470 & 6.3392 \\
2.5 & 223.96 & 0.001197 & 0.07995 & 962.0 & 2802.0 & 1840.1 & 2.5544 & 6.2560 \\
3 & 233.86 & 0.001217 & 0.06666 & 1008.4 & 2803.3 & 1794.9 & 2.6456 & 6.1858 \\
3.5 & 242.56 & 0.001235 & 0.05706 & 1049.8 & 2802.7 & 1753.0 & 2.7254 & 6.1245 \\
4 & 250.36 & 0.001253 & 0.04978 & 1087.4 & 2800.9 & 1713.5 & 2.7967 & 6.0697 \\
4.5 & 257.44 & 0.001270 & 0.04406 & 1122.1 & 2798.0 & 1675.9 & 2.8613 & 6.0198 \\
5 & 263.94 & 0.001286 & 0.03945 & 1154.5 & 2794.2 & 1639.7 & 2.9207 & 5.9737 \\
6 & 275.59 & 0.001319 & 0.03245 & 1213.7 & 2784.6 & 1570.8 & 3.0274 & 5.8901 \\
7 & 285.83 & 0.001352 & 0.02738 & 1267.4 & 2772.6 & 1505.1 & 3.1220 & 5.8146 \\
8 & 295.01 & 0.001385 & 0.02353 & 1317.1 & 2758.6 & 1441.5 & 3.2077 & 5.7448 \\
9 & 303.35 & 0.001418 & 0.02049 & 1363.7 & 2742.9 & 1379.2 & 3.2866 & 5.6790 \\
10 & 311.00 & 0.001453 & 0.01803 & 1407.9 & 2725.5 & 1317.6 & 3.3603 & 5.6159 \\
11 & 318.08 & 0.001489 & 0.01599 & 1450.3 & 2706.4 & 1256.1 & 3.4300 & 5.5545 \\
12 & 324.68 & 0.001526 & 0.01427 & 1491.3 & 2685.6 & 1194.3 & 3.4965 & 5.4941 \\
13 & 330.86 & 0.001566 & 0.01279 & 1531.4 & 2662.9 & 1131.5 & 3.5606 & 5.4339 \\
14 & 336.67 & 0.001610 & 0.01149 & 1570.9 & 2638.1 & 1067.2 & 3.6230 & 5.3730 \\
15 & 342.16 & 0.001657 & 0.01034 & 1610.2 & 2610.9 & 1000.7 & 3.6844 & 5.3108 \\
16 & 347.36 & 0.001710 & 0.009308 & 1649.7 & 2580.8 & 931.1 & 3.7457 & 5.2463 \\
17 & 352.29 & 0.001769 & 0.008369 & 1690.0 & 2547.4 & 857.4 & 3.8077 & 5.1785 \\
18 & 356.99 & 0.001839 & 0.007499 & 1732.0 & 2509.5 & 777.5 & 3.8717 & 5.1055 \\
19 & 361.47 & 0.001925 & 0.006673 & 1776.9 & 2465.4 & 688.5 & 3.9396 & 5.0246 \\
20 & 365.75 & 0.002039 & 0.005858 & 1827.1 & 2411.4 & 584.3 & 4.0154 & 4.9299 \\
21 & 369.83 & 0.002212 & 0.004988 & 1889.4 & 2337.5 & 448.1 & 4.1093 & 4.8062 \\
22.064\textsuperscript{*} & 373.95 & 0.003106 & 0.003106 & 2087.5 & 2087.5 & 0.0 & 4.4120 & 4.4120 \\
\end{longtblr}


\phbkhead{3　過熱蒸気表の抜き出し}
{\footnotesize\color{bkGray}色のついた行は、その圧力の飽和蒸気。「―」は、その温度では液（圧縮水）になる所。\qty{25}{MPa}は臨界圧力（\qty{22.064}{MPa}）より高いので、わく温度がない（熱-8 の超臨界圧の発電所）。表の間の値は、直線補間で読む（熱-8 ③）。\par}\vspace{1mm}
%% 自動生成（tools/mksteam2.py）。直接編集しないこと。値は IAPWS-IF97 の式（iapws 1.5.5）・CoolProp 8.0.0
\begin{longtblr}{colspec={Q[c,m,wd=10mm] Q[r,m] Q[r,m] Q[r,m] Q[r,m] Q[r,m] Q[r,m]}, width=\linewidth, rowhead=2, colsep=3.2pt, rowsep=0.3pt, hlines={0.3pt,bkLine}, vline{2,5}={0.4pt,bkLine}, cell{1}{2,5}={c=3}{c}, row{1-2}={bg=bkVol!14, font=\sffamily\bfseries\scriptsize, halign=c}, row{3}={bg=bkLine!35}, cells={font=\scriptsize}}
{温度\\$\theta$〔℃〕} & {$p=\qty{0.01}{MPa}$\\（飽和温度 \qty{45.81}{\celsius}）} & & & {$p=\qty{0.1}{MPa}$\\（飽和温度 \qty{99.61}{\celsius}）} & & \\
 & {$v$〔m$^3$/kg〕} & {$h$〔kJ/kg〕} & {$s$〔kJ/(kg·K)〕} & {$v$〔m$^3$/kg〕} & {$h$〔kJ/kg〕} & {$s$〔kJ/(kg·K)〕} \\
{\sffamily 飽和} & 14.67 & 2583.9 & 8.1489 & 1.694 & 2674.9 & 7.3588 \\
50 & 14.87 & 2592.0 & 8.1741 & ― & ― & ― \\
100 & 17.20 & 2687.4 & 8.4488 & ― & ― & ― \\
150 & 19.51 & 2783.0 & 8.6892 & 1.937 & 2776.6 & 7.6147 \\
200 & 21.83 & 2879.6 & 8.9048 & 2.172 & 2875.5 & 7.8356 \\
250 & 24.14 & 2977.4 & 9.1014 & 2.406 & 2974.5 & 8.0346 \\
300 & 26.45 & 3076.7 & 9.2827 & 2.639 & 3074.5 & 8.2171 \\
350 & 28.75 & 3177.5 & 9.4513 & 2.871 & 3175.8 & 8.3865 \\
400 & 31.06 & 3279.9 & 9.6093 & 3.103 & 3278.5 & 8.5451 \\
450 & 33.37 & 3384.0 & 9.7584 & 3.334 & 3382.8 & 8.6945 \\
500 & 35.68 & 3489.7 & 9.8997 & 3.566 & 3488.7 & 8.8361 \\
550 & 37.99 & 3597.1 & 10.0343 & 3.797 & 3596.3 & 8.9709 \\
600 & 40.30 & 3706.3 & 10.1631 & 4.028 & 3705.6 & 9.0998 \\
650 & 42.60 & 3817.2 & 10.2866 & 4.259 & 3816.6 & 9.2234 \\
700 & 44.91 & 3929.9 & 10.4055 & 4.490 & 3929.4 & 9.3424 \\
\end{longtblr}\par\vspace{1.5mm}
\begin{longtblr}{colspec={Q[c,m,wd=10mm] Q[r,m] Q[r,m] Q[r,m] Q[r,m] Q[r,m] Q[r,m]}, width=\linewidth, rowhead=2, colsep=3.2pt, rowsep=0.3pt, hlines={0.3pt,bkLine}, vline{2,5}={0.4pt,bkLine}, cell{1}{2,5}={c=3}{c}, row{1-2}={bg=bkVol!14, font=\sffamily\bfseries\scriptsize, halign=c}, row{3}={bg=bkLine!35}, cells={font=\scriptsize}}
{温度\\$\theta$〔℃〕} & {$p=\qty{0.2}{MPa}$\\（飽和温度 \qty{120.21}{\celsius}）} & & & {$p=\qty{0.5}{MPa}$\\（飽和温度 \qty{151.84}{\celsius}）} & & \\
 & {$v$〔m$^3$/kg〕} & {$h$〔kJ/kg〕} & {$s$〔kJ/(kg·K)〕} & {$v$〔m$^3$/kg〕} & {$h$〔kJ/kg〕} & {$s$〔kJ/(kg·K)〕} \\
{\sffamily 飽和} & 0.8857 & 2706.2 & 7.1269 & 0.3748 & 2748.1 & 6.8206 \\
150 & 0.9599 & 2769.1 & 7.2809 & ― & ― & ― \\
200 & 1.081 & 2870.8 & 7.5081 & 0.4250 & 2855.9 & 7.0611 \\
250 & 1.199 & 2971.3 & 7.7100 & 0.4744 & 2961.1 & 7.2726 \\
300 & 1.316 & 3072.1 & 7.8940 & 0.5226 & 3064.6 & 7.4614 \\
350 & 1.433 & 3173.9 & 8.0643 & 0.5701 & 3168.1 & 7.6345 \\
400 & 1.549 & 3277.0 & 8.2235 & 0.6173 & 3272.3 & 7.7954 \\
450 & 1.665 & 3381.5 & 8.3733 & 0.6642 & 3377.7 & 7.9464 \\
500 & 1.781 & 3487.6 & 8.5151 & 0.7109 & 3484.4 & 8.0891 \\
550 & 1.897 & 3595.4 & 8.6501 & 0.7576 & 3592.6 & 8.2247 \\
600 & 2.013 & 3704.8 & 8.7792 & 0.8041 & 3702.5 & 8.3543 \\
650 & 2.129 & 3815.9 & 8.9029 & 0.8506 & 3813.9 & 8.4784 \\
700 & 2.244 & 3928.8 & 9.0220 & 0.8970 & 3927.0 & 8.5977 \\
\end{longtblr}\par\vspace{1.5mm}
\begin{longtblr}{colspec={Q[c,m,wd=10mm] Q[r,m] Q[r,m] Q[r,m] Q[r,m] Q[r,m] Q[r,m]}, width=\linewidth, rowhead=2, colsep=3.2pt, rowsep=0.3pt, hlines={0.3pt,bkLine}, vline{2,5}={0.4pt,bkLine}, cell{1}{2,5}={c=3}{c}, row{1-2}={bg=bkVol!14, font=\sffamily\bfseries\scriptsize, halign=c}, row{3}={bg=bkLine!35}, cells={font=\scriptsize}}
{温度\\$\theta$〔℃〕} & {$p=\qty{1}{MPa}$\\（飽和温度 \qty{179.89}{\celsius}）} & & & {$p=\qty{2}{MPa}$\\（飽和温度 \qty{212.38}{\celsius}）} & & \\
 & {$v$〔m$^3$/kg〕} & {$h$〔kJ/kg〕} & {$s$〔kJ/(kg·K)〕} & {$v$〔m$^3$/kg〕} & {$h$〔kJ/kg〕} & {$s$〔kJ/(kg·K)〕} \\
{\sffamily 飽和} & 0.1943 & 2777.1 & 6.5850 & 0.09958 & 2798.4 & 6.3392 \\
200 & 0.2060 & 2828.3 & 6.6955 & ― & ― & ― \\
250 & 0.2327 & 2943.2 & 6.9266 & 0.1115 & 2903.2 & 6.5474 \\
300 & 0.2580 & 3051.7 & 7.1247 & 0.1255 & 3024.3 & 6.7685 \\
350 & 0.2825 & 3158.2 & 7.3028 & 0.1386 & 3137.6 & 6.9582 \\
400 & 0.3066 & 3264.4 & 7.4668 & 0.1512 & 3248.2 & 7.1290 \\
450 & 0.3304 & 3371.2 & 7.6198 & 0.1635 & 3358.1 & 7.2863 \\
500 & 0.3541 & 3479.0 & 7.7640 & 0.1757 & 3468.1 & 7.4335 \\
550 & 0.3777 & 3588.1 & 7.9007 & 0.1877 & 3578.9 & 7.5723 \\
600 & 0.4011 & 3698.6 & 8.0309 & 0.1996 & 3690.7 & 7.7042 \\
650 & 0.4245 & 3810.5 & 8.1557 & 0.2115 & 3803.8 & 7.8301 \\
700 & 0.4478 & 3924.1 & 8.2755 & 0.2233 & 3918.2 & 7.9509 \\
\end{longtblr}\par\vspace{1.5mm}
\begin{longtblr}{colspec={Q[c,m,wd=10mm] Q[r,m] Q[r,m] Q[r,m] Q[r,m] Q[r,m] Q[r,m]}, width=\linewidth, rowhead=2, colsep=3.2pt, rowsep=0.3pt, hlines={0.3pt,bkLine}, vline{2,5}={0.4pt,bkLine}, cell{1}{2,5}={c=3}{c}, row{1-2}={bg=bkVol!14, font=\sffamily\bfseries\scriptsize, halign=c}, row{3}={bg=bkLine!35}, cells={font=\scriptsize}}
{温度\\$\theta$〔℃〕} & {$p=\qty{4}{MPa}$\\（飽和温度 \qty{250.36}{\celsius}）} & & & {$p=\qty{10}{MPa}$\\（飽和温度 \qty{311.00}{\celsius}）} & & \\
 & {$v$〔m$^3$/kg〕} & {$h$〔kJ/kg〕} & {$s$〔kJ/(kg·K)〕} & {$v$〔m$^3$/kg〕} & {$h$〔kJ/kg〕} & {$s$〔kJ/(kg·K)〕} \\
{\sffamily 飽和} & 0.04978 & 2800.9 & 6.0697 & 0.01803 & 2725.5 & 5.6159 \\
300 & 0.05887 & 2961.7 & 6.3638 & ― & ― & ― \\
350 & 0.06647 & 3093.3 & 6.5843 & 0.02244 & 2924.0 & 5.9458 \\
400 & 0.07343 & 3214.4 & 6.7712 & 0.02644 & 3097.4 & 6.2139 \\
450 & 0.08004 & 3331.0 & 6.9383 & 0.02978 & 3242.3 & 6.4217 \\
500 & 0.08644 & 3445.8 & 7.0919 & 0.03281 & 3375.1 & 6.5993 \\
550 & 0.09270 & 3560.2 & 7.2353 & 0.03566 & 3501.9 & 6.7584 \\
600 & 0.09886 & 3674.8 & 7.3704 & 0.03838 & 3625.8 & 6.9045 \\
650 & 0.1049 & 3790.2 & 7.4989 & 0.04102 & 3748.3 & 7.0409 \\
700 & 0.1110 & 3906.4 & 7.6215 & 0.04359 & 3870.3 & 7.1696 \\
\end{longtblr}\par\vspace{1.5mm}
\begin{longtblr}{colspec={Q[c,m,wd=10mm] Q[r,m] Q[r,m] Q[r,m] Q[r,m] Q[r,m] Q[r,m]}, width=\linewidth, rowhead=2, colsep=3.2pt, rowsep=0.3pt, hlines={0.3pt,bkLine}, vline{2,5}={0.4pt,bkLine}, cell{1}{2,5}={c=3}{c}, row{1-2}={bg=bkVol!14, font=\sffamily\bfseries\scriptsize, halign=c}, row{3}={bg=bkLine!35}, cells={font=\scriptsize}}
{温度\\$\theta$〔℃〕} & {$p=\qty{15}{MPa}$\\（飽和温度 \qty{342.16}{\celsius}）} & & & {$p=\qty{25}{MPa}$\\（臨界圧力より高い）} & & \\
 & {$v$〔m$^3$/kg〕} & {$h$〔kJ/kg〕} & {$s$〔kJ/(kg·K)〕} & {$v$〔m$^3$/kg〕} & {$h$〔kJ/kg〕} & {$s$〔kJ/(kg·K)〕} \\
{\sffamily 飽和} & 0.01034 & 2610.9 & 5.3108 & ― & ― & ― \\
350 & 0.01148 & 2693.0 & 5.4435 & ― & ― & ― \\
400 & 0.01567 & 2975.5 & 5.8817 & 0.006005 & 2578.6 & 5.1399 \\
450 & 0.01848 & 3157.8 & 6.1433 & 0.009175 & 2950.4 & 5.6755 \\
500 & 0.02083 & 3310.8 & 6.3479 & 0.01114 & 3165.9 & 5.9642 \\
550 & 0.02295 & 3450.5 & 6.5230 & 0.01274 & 3339.3 & 6.1816 \\
600 & 0.02492 & 3583.3 & 6.6797 & 0.01414 & 3493.7 & 6.3638 \\
650 & 0.02680 & 3712.4 & 6.8235 & 0.01543 & 3638.0 & 6.5246 \\
700 & 0.02862 & 3839.5 & 6.9576 & 0.01664 & 3776.4 & 6.6706 \\
\end{longtblr}\par\vspace{1.5mm}


%% ---- 空気線図（見開き。左のページから始める）と湿り空気の表
%% 自動生成（tools/mksteam2.py）。直接編集しないこと。値は IAPWS-IF97 の式（iapws 1.5.5）・CoolProp 8.0.0
% 空気線図（101.325 kPa）。横：乾球温度、縦：絶対湿度。線：相対湿度・比エンタルピー・湿球温度・比体積（CoolProp の HAPropsSI）
\pgfplotsset{
  phPsyRH/.style={bkBlue!80!black, line width=0.6pt, mark=none},
  phPsySat/.style={bkBlue!80!black, line width=1.6pt, mark=none},
  phPsyH/.style={bkRed!75!black, line width=0.45pt, mark=none},
  phPsyWB/.style={bkGreen!70!black, line width=0.45pt, dash pattern=on 2.4pt off 1.6pt, mark=none},
  phPsyV/.style={bkPurple!75!black, line width=0.45pt, dash pattern=on 0.8pt off 1.2pt, mark=none},
  phPsyPt/.style={only marks, mark=*, mark size=1.5pt, bkInk, mark options={fill=white, line width=0.8pt}}}
\tikzset{phPsyRHlab/.style={font=\sffamily\scriptsize, text=bkBlue!80!black, inner sep=0.8pt},
  phPsyHlab/.style={font=\sffamily\scriptsize, text=bkRed!75!black, inner sep=0.6pt},
  phPsyVlab/.style={font=\sffamily\scriptsize, text=bkPurple!75!black, inner sep=0.8pt},
  phPsyWBlab/.style={font=\sffamily\scriptsize, text=bkGreen!50!black, inner sep=0.6pt},
  phPsyPtlab/.style={font=\sffamily\bfseries\scriptsize, text=bkInk, inner sep=0.6pt}}
\newcommand\phPsychroAxis[2]{% #1＝図のはば、#2＝高さ
\begin{tikzpicture}
\begin{axis}[name=psy, width=#1, height=#2, scale only axis, xmin=-10, xmax=50, ymin=0, ymax=30,
  axis lines=box, axis line style={bkInk, line width=0.7pt}, clip=true, clip mode=individual,
  xtick={-10,-5,...,50}, minor xtick={-9,...,49}, ytick={0,2,...,30}, minor ytick={1,3,...,29},
  yticklabel pos=left, ylabel near ticks,
  grid=both, major grid style={bkLine, line width=0.35pt}, minor grid style={bkLine!45, line width=0.2pt},
  tick label style={font=\sffamily\footnotesize, text=bkInk}, major tick length=2pt, minor tick length=1pt,
  xlabel={乾球温度 $\theta$〔℃〕}, xlabel style={at={(axis description cs:0.78,0)}, anchor=north, yshift=-9.5mm}, ylabel={絶対湿度 $x$〔g/kg(DA)〕},
  label style={font=\sffamily\small, text=bkInk}, /pgf/number format/1000 sep={}, scaled ticks=false]
\addplot[phPsyRH] coordinates {(-10.000,0.1602) (-9.500,0.1675) (-9.000,0.1751) (-8.500,0.1829) (-8.000,0.1911) (-7.500,0.1997) (-7.000,0.2085) (-6.500,0.2178) (-6.000,0.2274) (-5.500,0.2374) (-5.000,0.2477) (-4.500,0.2585) (-4.000,0.2698) (-3.500,0.2814) (-3.000,0.2936) (-2.500,0.3062) (-2.000,0.3193) (-1.500,0.3329) (-1.000,0.3470) (-0.500,0.3617) (0.000,0.3769) (0.500,0.3909) (1.000,0.4052) (1.500,0.4201) (2.000,0.4354) (2.500,0.4512) (3.000,0.4676) (3.500,0.4844) (4.000,0.5018) (4.500,0.5197) (5.000,0.5382) (5.500,0.5573) (6.000,0.5770) (6.500,0.5973) (7.000,0.6182) (7.500,0.6398) (8.000,0.6620) (8.500,0.6849) (9.000,0.7085) (9.500,0.7328) (10.000,0.7579) (10.500,0.7837) (11.000,0.8102) (11.500,0.8376) (12.000,0.8658) (12.500,0.8948) (13.000,0.9247) (13.500,0.9554) (14.000,0.9870) (14.500,1.0196) (15.000,1.0531) (15.500,1.0875) (16.000,1.1230) (16.500,1.1595) (17.000,1.1970) (17.500,1.2355) (18.000,1.2752) (18.500,1.3159) (19.000,1.3578) (19.500,1.4009) (20.000,1.4452) (20.500,1.4907) (21.000,1.5374) (21.500,1.5854) (22.000,1.6348) (22.500,1.6855) (23.000,1.7375) (23.500,1.7910) (24.000,1.8459) (24.500,1.9023) (25.000,1.9601) (25.500,2.0195) (26.000,2.0805) (26.500,2.1431) (27.000,2.2073) (27.500,2.2732) (28.000,2.3409) (28.500,2.4102) (29.000,2.4814) (29.500,2.5544) (30.000,2.6293) (30.250,2.6674) (30.500,2.7061) (30.750,2.7452) (31.000,2.7848) (31.250,2.8249) (31.500,2.8656) (31.750,2.9067) (32.000,2.9483) (32.250,2.9905) (32.500,3.0332) (32.750,3.0764) (33.000,3.1202) (33.250,3.1645) (33.500,3.2094) (33.750,3.2548) (34.000,3.3007) (34.250,3.3473) (34.500,3.3944) (34.750,3.4421) (35.000,3.4903) (35.250,3.5392) (35.500,3.5887) (35.750,3.6387) (36.000,3.6894) (36.250,3.7407) (36.500,3.7926) (36.750,3.8451) (37.000,3.8983) (37.250,3.9521) (37.500,4.0065) (37.750,4.0616) (38.000,4.1174) (38.250,4.1739) (38.500,4.2310) (38.750,4.2888) (39.000,4.3472) (39.250,4.4064) (39.500,4.4663) (39.750,4.5269) (40.000,4.5882) (40.250,4.6502) (40.500,4.7129) (40.750,4.7764) (41.000,4.8407) (41.250,4.9056) (41.500,4.9714) (41.750,5.0379) (42.000,5.1052) (42.250,5.1733) (42.500,5.2421) (42.750,5.3118) (43.000,5.3822) (43.250,5.4535) (43.500,5.5256) (43.750,5.5985) (44.000,5.6723) (44.250,5.7469) (44.500,5.8224) (44.750,5.8987) (45.000,5.9759) (45.250,6.0540) (45.500,6.1330) (45.750,6.2128) (46.000,6.2936) (46.250,6.3753) (46.500,6.4579) (46.750,6.5414) (47.000,6.6259) (47.250,6.7113) (47.500,6.7977) (47.750,6.8851) (48.000,6.9734) (48.250,7.0627) (48.500,7.1530) (48.750,7.2444) (49.000,7.3367) (49.250,7.4300) (49.500,7.5244) (49.750,7.6199) (50.000,7.7164)};
\addplot[phPsyRH] coordinates {(-10.000,0.3206) (-9.500,0.3351) (-9.000,0.3502) (-8.500,0.3660) (-8.000,0.3824) (-7.500,0.3994) (-7.000,0.4172) (-6.500,0.4357) (-6.000,0.4549) (-5.500,0.4749) (-5.000,0.4957) (-4.500,0.5173) (-4.000,0.5398) (-3.500,0.5631) (-3.000,0.5874) (-2.500,0.6127) (-2.000,0.6389) (-1.500,0.6661) (-1.000,0.6944) (-0.500,0.7238) (0.000,0.7543) (0.500,0.7822) (1.000,0.8110) (1.500,0.8407) (2.000,0.8714) (2.500,0.9031) (3.000,0.9358) (3.500,0.9696) (4.000,1.0044) (4.500,1.0403) (5.000,1.0774) (5.500,1.1156) (6.000,1.1551) (6.500,1.1957) (7.000,1.2376) (7.500,1.2809) (8.000,1.3254) (8.500,1.3713) (9.000,1.4186) (9.500,1.4674) (10.000,1.5176) (10.500,1.5693) (11.000,1.6226) (11.500,1.6775) (12.000,1.7340) (12.500,1.7922) (13.000,1.8521) (13.500,1.9137) (14.000,1.9772) (14.500,2.0425) (15.000,2.1097) (15.500,2.1789) (16.000,2.2500) (16.500,2.3232) (17.000,2.3985) (17.500,2.4760) (18.000,2.5556) (18.500,2.6374) (19.000,2.7216) (19.500,2.8081) (20.000,2.8971) (20.500,2.9885) (21.000,3.0825) (21.500,3.1790) (22.000,3.2782) (22.500,3.3801) (23.000,3.4848) (23.500,3.5923) (24.000,3.7028) (24.500,3.8162) (25.000,3.9326) (25.500,4.0522) (26.000,4.1750) (26.500,4.3010) (27.000,4.4304) (27.500,4.5631) (28.000,4.6994) (28.500,4.8392) (29.000,4.9827) (29.500,5.1299) (30.000,5.2809) (30.250,5.3579) (30.500,5.4358) (30.750,5.5148) (31.000,5.5947) (31.250,5.6757) (31.500,5.7577) (31.750,5.8407) (32.000,5.9248) (32.250,6.0099) (32.500,6.0961) (32.750,6.1835) (33.000,6.2719) (33.250,6.3614) (33.500,6.4520) (33.750,6.5438) (34.000,6.6367) (34.250,6.7308) (34.500,6.8260) (34.750,6.9225) (35.000,7.0201) (35.250,7.1189) (35.500,7.2190) (35.750,7.3203) (36.000,7.4228) (36.250,7.5266) (36.500,7.6317) (36.750,7.7381) (37.000,7.8457) (37.250,7.9547) (37.500,8.0650) (37.750,8.1767) (38.000,8.2897) (38.250,8.4041) (38.500,8.5199) (38.750,8.6371) (39.000,8.7557) (39.250,8.8757) (39.500,8.9972) (39.750,9.1201) (40.000,9.2445) (40.250,9.3704) (40.500,9.4978) (40.750,9.6268) (41.000,9.7573) (41.250,9.8893) (41.500,10.0229) (41.750,10.1581) (42.000,10.2949) (42.250,10.4333) (42.500,10.5733) (42.750,10.7151) (43.000,10.8584) (43.250,11.0035) (43.500,11.1503) (43.750,11.2988) (44.000,11.4490) (44.250,11.6010) (44.500,11.7548) (44.750,11.9104) (45.000,12.0678) (45.250,12.2270) (45.500,12.3881) (45.750,12.5510) (46.000,12.7159) (46.250,12.8826) (46.500,13.0513) (46.750,13.2219) (47.000,13.3945) (47.250,13.5691) (47.500,13.7457) (47.750,13.9243) (48.000,14.1050) (48.250,14.2877) (48.500,14.4725) (48.750,14.6595) (49.000,14.8485) (49.250,15.0398) (49.500,15.2332) (49.750,15.4288) (50.000,15.6266)};
\addplot[phPsyRH] coordinates {(-10.000,0.4810) (-9.500,0.5028) (-9.000,0.5255) (-8.500,0.5491) (-8.000,0.5737) (-7.500,0.5994) (-7.000,0.6260) (-6.500,0.6537) (-6.000,0.6826) (-5.500,0.7126) (-5.000,0.7438) (-4.500,0.7763) (-4.000,0.8100) (-3.500,0.8451) (-3.000,0.8815) (-2.500,0.9194) (-2.000,0.9588) (-1.500,0.9997) (-1.000,1.0422) (-0.500,1.0863) (0.000,1.1322) (0.500,1.1741) (1.000,1.2173) (1.500,1.2620) (2.000,1.3081) (2.500,1.3557) (3.000,1.4048) (3.500,1.4555) (4.000,1.5078) (4.500,1.5618) (5.000,1.6175) (5.500,1.6750) (6.000,1.7342) (6.500,1.7953) (7.000,1.8583) (7.500,1.9233) (8.000,1.9902) (8.500,2.0592) (9.000,2.1303) (9.500,2.2036) (10.000,2.2791) (10.500,2.3569) (11.000,2.4371) (11.500,2.5196) (12.000,2.6046) (12.500,2.6921) (13.000,2.7822) (13.500,2.8750) (14.000,2.9705) (14.500,3.0688) (15.000,3.1700) (15.500,3.2741) (16.000,3.3812) (16.500,3.4914) (17.000,3.6047) (17.500,3.7213) (18.000,3.8413) (18.500,3.9646) (19.000,4.0914) (19.500,4.2218) (20.000,4.3558) (20.500,4.4936) (21.000,4.6352) (21.500,4.7807) (22.000,4.9303) (22.500,5.0840) (23.000,5.2419) (23.500,5.4041) (24.000,5.5707) (24.500,5.7419) (25.000,5.9177) (25.500,6.0982) (26.000,6.2836) (26.500,6.4739) (27.000,6.6693) (27.500,6.8699) (28.000,7.0758) (28.500,7.2872) (29.000,7.5041) (29.500,7.7267) (30.000,7.9551) (30.250,8.0716) (30.500,8.1895) (30.750,8.3090) (31.000,8.4300) (31.250,8.5525) (31.500,8.6767) (31.750,8.8024) (32.000,8.9297) (32.250,9.0587) (32.500,9.1893) (32.750,9.3215) (33.000,9.4555) (33.250,9.5911) (33.500,9.7285) (33.750,9.8676) (34.000,10.0085) (34.250,10.1511) (34.500,10.2955) (34.750,10.4418) (35.000,10.5899) (35.250,10.7398) (35.500,10.8917) (35.750,11.0454) (36.000,11.2011) (36.250,11.3587) (36.500,11.5182) (36.750,11.6798) (37.000,11.8433) (37.250,12.0089) (37.500,12.1765) (37.750,12.3462) (38.000,12.5180) (38.250,12.6919) (38.500,12.8680) (38.750,13.0462) (39.000,13.2266) (39.250,13.4092) (39.500,13.5941) (39.750,13.7812) (40.000,13.9706) (40.250,14.1623) (40.500,14.3564) (40.750,14.5528) (41.000,14.7516) (41.250,14.9528) (41.500,15.1565) (41.750,15.3626) (42.000,15.5712) (42.250,15.7823) (42.500,15.9960) (42.750,16.2122) (43.000,16.4311) (43.250,16.6526) (43.500,16.8767) (43.750,17.1035) (44.000,17.3331) (44.250,17.5654) (44.500,17.8004) (44.750,18.0383) (45.000,18.2790) (45.250,18.5226) (45.500,18.7691) (45.750,19.0185) (46.000,19.2708) (46.250,19.5262) (46.500,19.7845) (46.750,20.0460) (47.000,20.3105) (47.250,20.5781) (47.500,20.8489) (47.750,21.1229) (48.000,21.4001) (48.250,21.6806) (48.500,21.9644) (48.750,22.2514) (49.000,22.5419) (49.250,22.8358) (49.500,23.1331) (49.750,23.4338) (50.000,23.7381)};
\addplot[phPsyRH] coordinates {(-10.000,0.6415) (-9.500,0.6706) (-9.000,0.7009) (-8.500,0.7324) (-8.000,0.7652) (-7.500,0.7994) (-7.000,0.8350) (-6.500,0.8720) (-6.000,0.9105) (-5.500,0.9505) (-5.000,0.9921) (-4.500,1.0355) (-4.000,1.0805) (-3.500,1.1273) (-3.000,1.1760) (-2.500,1.2265) (-2.000,1.2791) (-1.500,1.3337) (-1.000,1.3904) (-0.500,1.4493) (0.000,1.5105) (0.500,1.5664) (1.000,1.6241) (1.500,1.6838) (2.000,1.7453) (2.500,1.8089) (3.000,1.8745) (3.500,1.9422) (4.000,2.0121) (4.500,2.0842) (5.000,2.1586) (5.500,2.2353) (6.000,2.3144) (6.500,2.3961) (7.000,2.4802) (7.500,2.5670) (8.000,2.6564) (8.500,2.7487) (9.000,2.8437) (9.500,2.9416) (10.000,3.0426) (10.500,3.1465) (11.000,3.2537) (11.500,3.3640) (12.000,3.4777) (12.500,3.5947) (13.000,3.7152) (13.500,3.8393) (14.000,3.9670) (14.500,4.0985) (15.000,4.2338) (15.500,4.3731) (16.000,4.5164) (16.500,4.6639) (17.000,4.8156) (17.500,4.9717) (18.000,5.1322) (18.500,5.2974) (19.000,5.4672) (19.500,5.6418) (20.000,5.8213) (20.500,6.0059) (21.000,6.1956) (21.500,6.3907) (22.000,6.5911) (22.500,6.7971) (23.000,7.0089) (23.500,7.2264) (24.000,7.4499) (24.500,7.6795) (25.000,7.9153) (25.500,8.1576) (26.000,8.4064) (26.500,8.6619) (27.000,8.9243) (27.500,9.1937) (28.000,9.4704) (28.500,9.7543) (29.000,10.0459) (29.500,10.3451) (30.000,10.6523) (30.250,10.8089) (30.500,10.9675) (30.750,11.1282) (31.000,11.2910) (31.250,11.4559) (31.500,11.6229) (31.750,11.7921) (32.000,11.9635) (32.250,12.1371) (32.500,12.3130) (32.750,12.4911) (33.000,12.6715) (33.250,12.8542) (33.500,13.0393) (33.750,13.2267) (34.000,13.4166) (34.250,13.6088) (34.500,13.8036) (34.750,14.0007) (35.000,14.2005) (35.250,14.4027) (35.500,14.6075) (35.750,14.8149) (36.000,15.0249) (36.250,15.2376) (36.500,15.4530) (36.750,15.6711) (37.000,15.8919) (37.250,16.1156) (37.500,16.3420) (37.750,16.5713) (38.000,16.8034) (38.250,17.0384) (38.500,17.2764) (38.750,17.5174) (39.000,17.7614) (39.250,18.0084) (39.500,18.2585) (39.750,18.5117) (40.000,18.7680) (40.250,19.0275) (40.500,19.2903) (40.750,19.5562) (41.000,19.8255) (41.250,20.0981) (41.500,20.3741) (41.750,20.6535) (42.000,20.9363) (42.250,21.2226) (42.500,21.5124) (42.750,21.8058) (43.000,22.1028) (43.250,22.4034) (43.500,22.7077) (43.750,23.0157) (44.000,23.3275) (44.250,23.6431) (44.500,23.9625) (44.750,24.2859) (45.000,24.6132) (45.250,24.9444) (45.500,25.2797) (45.750,25.6191) (46.000,25.9626) (46.250,26.3102) (46.500,26.6621) (46.750,27.0182) (47.000,27.3787) (47.250,27.7435) (47.500,28.1127) (47.750,28.4864) (48.000,28.8646) (48.250,29.2473) (48.500,29.6347) (48.733,30.0000)};
\addplot[phPsyRH] coordinates {(-10.000,0.8021) (-9.500,0.8384) (-9.000,0.8763) (-8.500,0.9158) (-8.000,0.9568) (-7.500,0.9996) (-7.000,1.0440) (-6.500,1.0903) (-6.000,1.1385) (-5.500,1.1886) (-5.000,1.2407) (-4.500,1.2949) (-4.000,1.3512) (-3.500,1.4098) (-3.000,1.4706) (-2.500,1.5339) (-2.000,1.5997) (-1.500,1.6680) (-1.000,1.7390) (-0.500,1.8127) (0.000,1.8893) (0.500,1.9592) (1.000,2.0315) (1.500,2.1061) (2.000,2.1832) (2.500,2.2628) (3.000,2.3449) (3.500,2.4297) (4.000,2.5171) (4.500,2.6074) (5.000,2.7005) (5.500,2.7966) (6.000,2.8957) (6.500,2.9980) (7.000,3.1034) (7.500,3.2121) (8.000,3.3241) (8.500,3.4396) (9.000,3.5587) (9.500,3.6814) (10.000,3.8079) (10.500,3.9382) (11.000,4.0724) (11.500,4.2107) (12.000,4.3531) (12.500,4.4999) (13.000,4.6509) (13.500,4.8065) (14.000,4.9667) (14.500,5.1316) (15.000,5.3013) (15.500,5.4760) (16.000,5.6558) (16.500,5.8408) (17.000,6.0312) (17.500,6.2271) (18.000,6.4286) (18.500,6.6358) (19.000,6.8490) (19.500,7.0682) (20.000,7.2937) (20.500,7.5255) (21.000,7.7639) (21.500,8.0089) (22.000,8.2608) (22.500,8.5197) (23.000,8.7858) (23.500,9.0593) (24.000,9.3403) (24.500,9.6291) (25.000,9.9257) (25.500,10.2305) (26.000,10.5436) (26.500,10.8652) (27.000,11.1956) (27.500,11.5348) (28.000,11.8832) (28.500,12.2409) (29.000,12.6083) (29.500,12.9854) (30.000,13.3726) (30.250,13.5700) (30.500,13.7701) (30.750,13.9727) (31.000,14.1781) (31.250,14.3861) (31.500,14.5968) (31.750,14.8104) (32.000,15.0267) (32.250,15.2458) (32.500,15.4678) (32.750,15.6927) (33.000,15.9205) (33.250,16.1512) (33.500,16.3850) (33.750,16.6218) (34.000,16.8617) (34.250,17.1046) (34.500,17.3507) (34.750,17.6000) (35.000,17.8525) (35.250,18.1082) (35.500,18.3672) (35.750,18.6296) (36.000,18.8953) (36.250,19.1644) (36.500,19.4370) (36.750,19.7131) (37.000,19.9926) (37.250,20.2758) (37.500,20.5626) (37.750,20.8530) (38.000,21.1471) (38.250,21.4449) (38.500,21.7466) (38.750,22.0520) (39.000,22.3614) (39.250,22.6746) (39.500,22.9918) (39.750,23.3130) (40.000,23.6383) (40.250,23.9677) (40.500,24.3013) (40.750,24.6390) (41.000,24.9810) (41.250,25.3273) (41.500,25.6779) (41.750,26.0330) (42.000,26.3925) (42.250,26.7565) (42.500,27.1251) (42.750,27.4983) (43.000,27.8761) (43.250,28.2587) (43.500,28.6461) (43.750,29.0383) (44.000,29.4354) (44.250,29.8374) (44.350,30.0000)};
\addplot[phPsyRH] coordinates {(-10.000,0.9627) (-9.500,1.0064) (-9.000,1.0519) (-8.500,1.0992) (-8.000,1.1485) (-7.500,1.1999) (-7.000,1.2533) (-6.500,1.3089) (-6.000,1.3667) (-5.500,1.4268) (-5.000,1.4894) (-4.500,1.5545) (-4.000,1.6221) (-3.500,1.6925) (-3.000,1.7656) (-2.500,1.8416) (-2.000,1.9206) (-1.500,2.0027) (-1.000,2.0879) (-0.500,2.1765) (0.000,2.2685) (0.500,2.3526) (1.000,2.4394) (1.500,2.5291) (2.000,2.6217) (2.500,2.7173) (3.000,2.8160) (3.500,2.9179) (4.000,3.0230) (4.500,3.1315) (5.000,3.2435) (5.500,3.3590) (6.000,3.4781) (6.500,3.6010) (7.000,3.7278) (7.500,3.8585) (8.000,3.9932) (8.500,4.1321) (9.000,4.2753) (9.500,4.4229) (10.000,4.5750) (10.500,4.7318) (11.000,4.8933) (11.500,5.0597) (12.000,5.2311) (12.500,5.4077) (13.000,5.5895) (13.500,5.7767) (14.000,5.9695) (14.500,6.1681) (15.000,6.3724) (15.500,6.5828) (16.000,6.7993) (16.500,7.0222) (17.000,7.2515) (17.500,7.4875) (18.000,7.7302) (18.500,7.9800) (19.000,8.2369) (19.500,8.5012) (20.000,8.7730) (20.500,9.0525) (21.000,9.3400) (21.500,9.6355) (22.000,9.9394) (22.500,10.2517) (23.000,10.5728) (23.500,10.9029) (24.000,11.2421) (24.500,11.5908) (25.000,11.9490) (25.500,12.3172) (26.000,12.6954) (26.500,13.0840) (27.000,13.4832) (27.500,13.8933) (28.000,14.3145) (28.500,14.7472) (29.000,15.1915) (29.500,15.6478) (30.000,16.1164) (30.250,16.3554) (30.500,16.5976) (30.750,16.8430) (31.000,17.0916) (31.250,17.3435) (31.500,17.5988) (31.750,17.8575) (32.000,18.1196) (32.250,18.3851) (32.500,18.6541) (32.750,18.9267) (33.000,19.2029) (33.250,19.4827) (33.500,19.7661) (33.750,20.0533) (34.000,20.3443) (34.250,20.6391) (34.500,20.9377) (34.750,21.2402) (35.000,21.5467) (35.250,21.8571) (35.500,22.1716) (35.750,22.4902) (36.000,22.8130) (36.250,23.1399) (36.500,23.4711) (36.750,23.8066) (37.000,24.1464) (37.250,24.4906) (37.500,24.8393) (37.750,25.1925) (38.000,25.5502) (38.250,25.9126) (38.500,26.2797) (38.750,26.6514) (39.000,27.0280) (39.250,27.4094) (39.500,27.7957) (39.750,28.1870) (40.000,28.5833) (40.250,28.9847) (40.500,29.3912) (40.750,29.8029) (40.869,30.0000)};
\addplot[phPsyRH] coordinates {(-10.000,1.1235) (-9.500,1.1745) (-9.000,1.2275) (-8.500,1.2828) (-8.000,1.3404) (-7.500,1.4003) (-7.000,1.4627) (-6.500,1.5275) (-6.000,1.5950) (-5.500,1.6653) (-5.000,1.7383) (-4.500,1.8143) (-4.000,1.8933) (-3.500,1.9755) (-3.000,2.0608) (-2.500,2.1496) (-2.000,2.2418) (-1.500,2.3377) (-1.000,2.4373) (-0.500,2.5407) (0.000,2.6482) (0.500,2.7464) (1.000,2.8478) (1.500,2.9526) (2.000,3.0608) (2.500,3.1725) (3.000,3.2878) (3.500,3.4068) (4.000,3.5297) (4.500,3.6565) (5.000,3.7873) (5.500,3.9223) (6.000,4.0616) (6.500,4.2052) (7.000,4.3534) (7.500,4.5062) (8.000,4.6637) (8.500,4.8262) (9.000,4.9936) (9.500,5.1662) (10.000,5.3441) (10.500,5.5274) (11.000,5.7163) (11.500,5.9110) (12.000,6.1115) (12.500,6.3181) (13.000,6.5309) (13.500,6.7500) (14.000,6.9756) (14.500,7.2080) (15.000,7.4472) (15.500,7.6935) (16.000,7.9470) (16.500,8.2080) (17.000,8.4766) (17.500,8.7529) (18.000,9.0373) (18.500,9.3300) (19.000,9.6310) (19.500,9.9407) (20.000,10.2593) (20.500,10.5870) (21.000,10.9240) (21.500,11.2705) (22.000,11.6269) (22.500,11.9933) (23.000,12.3700) (23.500,12.7573) (24.000,13.1555) (24.500,13.5647) (25.000,13.9853) (25.500,14.4176) (26.000,14.8619) (26.500,15.3184) (27.000,15.7875) (27.500,16.2694) (28.000,16.7646) (28.500,17.2733) (29.000,17.7959) (29.500,18.3327) (30.000,18.8840) (30.250,19.1653) (30.500,19.4503) (30.750,19.7392) (31.000,20.0319) (31.250,20.3286) (31.500,20.6292) (31.750,20.9339) (32.000,21.2426) (32.250,21.5555) (32.500,21.8725) (32.750,22.1937) (33.000,22.5192) (33.250,22.8491) (33.500,23.1833) (33.750,23.5220) (34.000,23.8651) (34.250,24.2128) (34.500,24.5651) (34.750,24.9221) (35.000,25.2838) (35.250,25.6502) (35.500,26.0215) (35.750,26.3977) (36.000,26.7789) (36.250,27.1650) (36.500,27.5563) (36.750,27.9527) (37.000,28.3543) (37.250,28.7612) (37.500,29.1734) (37.750,29.5910) (37.992,30.0000)};
\addplot[phPsyRH] coordinates {(-10.000,1.2843) (-9.500,1.3426) (-9.000,1.4033) (-8.500,1.4665) (-8.000,1.5323) (-7.500,1.6009) (-7.000,1.6722) (-6.500,1.7464) (-6.000,1.8236) (-5.500,1.9039) (-5.000,1.9875) (-4.500,2.0744) (-4.000,2.1647) (-3.500,2.2587) (-3.000,2.3564) (-2.500,2.4579) (-2.000,2.5634) (-1.500,2.6731) (-1.000,2.7870) (-0.500,2.9054) (0.000,3.0283) (0.500,3.1407) (1.000,3.2568) (1.500,3.3767) (2.000,3.5005) (2.500,3.6283) (3.000,3.7603) (3.500,3.8966) (4.000,4.0372) (4.500,4.1824) (5.000,4.3322) (5.500,4.4867) (6.000,4.6462) (6.500,4.8106) (7.000,4.9803) (7.500,5.1553) (8.000,5.3357) (8.500,5.5217) (9.000,5.7135) (9.500,5.9112) (10.000,6.1151) (10.500,6.3251) (11.000,6.5416) (11.500,6.7646) (12.000,6.9944) (12.500,7.2312) (13.000,7.4750) (13.500,7.7262) (14.000,7.9849) (14.500,8.2514) (15.000,8.5257) (15.500,8.8081) (16.000,9.0989) (16.500,9.3983) (17.000,9.7064) (17.500,10.0235) (18.000,10.3499) (18.500,10.6857) (19.000,11.0313) (19.500,11.3868) (20.000,11.7526) (20.500,12.1289) (21.000,12.5159) (21.500,12.9140) (22.000,13.3235) (22.500,13.7445) (23.000,14.1775) (23.500,14.6227) (24.000,15.0804) (24.500,15.5510) (25.000,16.0347) (25.500,16.5320) (26.000,17.0432) (26.500,17.5685) (27.000,18.1085) (27.500,18.6634) (28.000,19.2336) (28.500,19.8195) (29.000,20.4216) (29.500,21.0402) (30.000,21.6758) (30.250,22.0000) (30.500,22.3287) (30.750,22.6618) (31.000,22.9995) (31.250,23.3417) (31.500,23.6885) (31.750,24.0401) (32.000,24.3963) (32.250,24.7574) (32.500,25.1233) (32.750,25.4942) (33.000,25.8701) (33.250,26.2510) (33.500,26.6370) (33.750,27.0283) (34.000,27.4248) (34.250,27.8265) (34.500,28.2337) (34.750,28.6464) (35.000,29.0645) (35.250,29.4883) (35.500,29.9177) (35.548,30.0000)};
\addplot[phPsyRH] coordinates {(-10.000,1.4452) (-9.500,1.5108) (-9.000,1.5792) (-8.500,1.6503) (-8.000,1.7244) (-7.500,1.8015) (-7.000,1.8818) (-6.500,1.9654) (-6.000,2.0523) (-5.500,2.1427) (-5.000,2.2368) (-4.500,2.3346) (-4.000,2.4364) (-3.500,2.5422) (-3.000,2.6522) (-2.500,2.7665) (-2.000,2.8853) (-1.500,3.0088) (-1.000,3.1371) (-0.500,3.2705) (0.000,3.4090) (0.500,3.5355) (1.000,3.6663) (1.500,3.8014) (2.000,3.9408) (2.500,4.0849) (3.000,4.2336) (3.500,4.3871) (4.000,4.5456) (4.500,4.7091) (5.000,4.8779) (5.500,5.0521) (6.000,5.2318) (6.500,5.4172) (7.000,5.6084) (7.500,5.8057) (8.000,6.0091) (8.500,6.2188) (9.000,6.4351) (9.500,6.6581) (10.000,6.8879) (10.500,7.1248) (11.000,7.3689) (11.500,7.6205) (12.000,7.8798) (12.500,8.1469) (13.000,8.4221) (13.500,8.7055) (14.000,8.9975) (14.500,9.2982) (15.000,9.6079) (15.500,9.9267) (16.000,10.2551) (16.500,10.5931) (17.000,10.9410) (17.500,11.2992) (18.000,11.6679) (18.500,12.0473) (19.000,12.4378) (19.500,12.8396) (20.000,13.2530) (20.500,13.6783) (21.000,14.1159) (21.500,14.5661) (22.000,15.0291) (22.500,15.5054) (23.000,15.9952) (23.500,16.4990) (24.000,17.0170) (24.500,17.5497) (25.000,18.0974) (25.500,18.6605) (26.000,19.2395) (26.500,19.8346) (27.000,20.4464) (27.500,21.0753) (28.000,21.7218) (28.500,22.3862) (29.000,23.0690) (29.500,23.7708) (30.000,24.4919) (30.250,24.8600) (30.500,25.2330) (30.750,25.6112) (31.000,25.9946) (31.250,26.3832) (31.500,26.7771) (31.750,27.1764) (32.000,27.5811) (32.250,27.9914) (32.500,28.4072) (32.750,28.8287) (33.000,29.2560) (33.250,29.6890) (33.427,30.0000)};
\addplot[phPsySat] coordinates {(-10.000,1.6062) (-9.500,1.6791) (-9.000,1.7551) (-8.500,1.8342) (-8.000,1.9166) (-7.500,2.0024) (-7.000,2.0916) (-6.500,2.1845) (-6.000,2.2811) (-5.500,2.3817) (-5.000,2.4863) (-4.500,2.5951) (-4.000,2.7083) (-3.500,2.8259) (-3.000,2.9482) (-2.500,3.0754) (-2.000,3.2076) (-1.500,3.3449) (-1.000,3.4877) (-0.500,3.6360) (0.000,3.7900) (0.500,3.9309) (1.000,4.0763) (1.500,4.2266) (2.000,4.3818) (2.500,4.5421) (3.000,4.7075) (3.500,4.8784) (4.000,5.0547) (4.500,5.2368) (5.000,5.4247) (5.500,5.6185) (6.000,5.8186) (6.500,6.0250) (7.000,6.2379) (7.500,6.4575) (8.000,6.6839) (8.500,6.9175) (9.000,7.1584) (9.500,7.4067) (10.000,7.6626) (10.500,7.9265) (11.000,8.1985) (11.500,8.4788) (12.000,8.7677) (12.500,9.0653) (13.000,9.3720) (13.500,9.6879) (14.000,10.0133) (14.500,10.3485) (15.000,10.6938) (15.500,11.0493) (16.000,11.4154) (16.500,11.7924) (17.000,12.1805) (17.500,12.5801) (18.000,12.9914) (18.500,13.4148) (19.000,13.8505) (19.500,14.2990) (20.000,14.7605) (20.500,15.2354) (21.000,15.7240) (21.500,16.2268) (22.000,16.7440) (22.500,17.2761) (23.000,17.8234) (23.500,18.3864) (24.000,18.9654) (24.500,19.5610) (25.000,20.1734) (25.500,20.8033) (26.000,21.4509) (26.500,22.1168) (27.000,22.8016) (27.500,23.5055) (28.000,24.2293) (28.500,24.9734) (29.000,25.7383) (29.500,26.5246) (30.000,27.3329) (30.250,27.7454) (30.500,28.1637) (30.750,28.5877) (31.000,29.0176) (31.250,29.4535) (31.500,29.8953) (31.559,30.0000)};
\addplot[phPsyH] coordinates {(-5.763,2.3283) (-5.500,2.2214) (-5.000,2.0187) (-4.500,1.8161) (-4.000,1.6137) (-3.500,1.4115) (-3.000,1.2094) (-2.500,1.0074) (-2.000,0.8056) (-1.500,0.6040) (-1.000,0.4025) (-0.500,0.2012) (0.000,0.0000)};
\addplot[phPsyH] coordinates {(0.307,3.8759) (0.500,3.7976) (1.000,3.5951) (1.500,3.3928) (2.000,3.1907) (2.500,2.9887) (3.000,2.7869) (3.500,2.5852) (4.000,2.3837) (4.500,2.1823) (5.000,1.9811) (5.500,1.7800) (6.000,1.5791) (6.500,1.3783) (7.000,1.1777) (7.500,0.9772) (8.000,0.7769) (8.500,0.5767) (9.000,0.3767) (9.500,0.1769)};
\addplot[phPsyH] coordinates {(5.682,5.6906) (6.000,5.5618) (6.500,5.3595) (7.000,5.1573) (7.500,4.9553) (8.000,4.7534) (8.500,4.5517) (9.000,4.3502) (9.500,4.1488) (10.000,3.9475) (10.500,3.7464) (11.000,3.5454) (11.500,3.3446) (12.000,3.1440) (12.500,2.9435) (13.000,2.7431) (13.500,2.5429) (14.000,2.3429) (14.500,2.1430) (15.000,1.9432) (15.500,1.7436) (16.000,1.5441) (16.500,1.3448) (17.000,1.1457) (17.500,0.9466) (18.000,0.7478) (18.500,0.5490) (19.000,0.3505) (19.500,0.1520)};
\addplot[phPsyH] coordinates {(10.276,7.8072) (10.500,7.7164) (11.000,7.5139) (11.500,7.3115) (12.000,7.1093) (12.500,6.9073) (13.000,6.7054) (13.500,6.5036) (14.000,6.3021) (14.500,6.1006) (15.000,5.8993) (15.500,5.6982) (16.000,5.4972) (16.500,5.2964) (17.000,5.0957) (17.500,4.8952) (18.000,4.6948) (18.500,4.4945) (19.000,4.2944) (19.500,4.0945) (20.000,3.8947) (20.500,3.6951) (21.000,3.4956) (21.500,3.2962) (22.000,3.0970) (22.500,2.8980) (23.000,2.6990) (23.500,2.5003) (24.000,2.3017) (24.500,2.1032) (25.000,1.9049) (25.500,1.7067) (26.000,1.5086) (26.500,1.3107) (27.000,1.1130) (27.500,0.9154) (28.000,0.7179) (28.500,0.5206) (29.000,0.3234) (29.500,0.1264)};
\addplot[phPsyH] coordinates {(14.232,10.1679) (14.500,10.0594) (15.000,9.8565) (15.500,9.6538) (16.000,9.4513) (16.500,9.2490) (17.000,9.0467) (17.500,8.8447) (18.000,8.6428) (18.500,8.4410) (19.000,8.2394) (19.500,8.0379) (20.000,7.8366) (20.500,7.6355) (21.000,7.4344) (21.500,7.2336) (22.000,7.0329) (22.500,6.8323) (23.000,6.6319) (23.500,6.4316) (24.000,6.2315) (24.500,6.0315) (25.000,5.8317) (25.500,5.6320) (26.000,5.4325) (26.500,5.2331) (27.000,5.0339) (27.500,4.8348) (28.000,4.6358) (28.500,4.4370) (29.000,4.2384) (29.500,4.0399) (30.000,3.8415) (30.500,3.6433) (31.000,3.4452) (31.500,3.2473) (32.000,3.0495) (32.500,2.8518) (33.000,2.6543) (33.500,2.4569) (34.000,2.2597) (34.500,2.0626) (35.000,1.8657) (35.500,1.6689) (36.000,1.4722) (36.500,1.2757) (37.000,1.0793) (37.500,0.8831) (38.000,0.6870) (38.500,0.4911) (39.000,0.2952) (39.500,0.0996)};
\addplot[phPsyH] coordinates {(17.676,12.7235) (18.000,12.5917) (18.500,12.3884) (19.000,12.1853) (19.500,11.9823) (20.000,11.7794) (20.500,11.5768) (21.000,11.3742) (21.500,11.1718) (22.000,10.9696) (22.500,10.7675) (23.000,10.5656) (23.500,10.3638) (24.000,10.1622) (24.500,9.9607) (25.000,9.7594) (25.500,9.5582) (26.000,9.3572) (26.500,9.1563) (27.000,8.9556) (27.500,8.7550) (28.000,8.5546) (28.500,8.3543) (29.000,8.1541) (29.500,7.9541) (30.000,7.7543) (30.500,7.5546) (31.000,7.3550) (31.500,7.1556) (32.000,6.9563) (32.500,6.7572) (33.000,6.5582) (33.500,6.3594) (34.000,6.1607) (34.500,5.9622) (35.000,5.7638) (35.500,5.5655) (36.000,5.3674) (36.500,5.1694) (37.000,4.9716) (37.500,4.7739) (38.000,4.5763) (38.500,4.3789) (39.000,4.1817) (39.500,3.9845) (40.000,3.7876) (40.500,3.5907) (41.000,3.3940) (41.500,3.1975) (42.000,3.0010) (42.500,2.8048) (43.000,2.6086) (43.500,2.4126) (44.000,2.2167) (44.500,2.0210) (45.000,1.8254) (45.500,1.6300) (46.000,1.4347) (46.500,1.2395) (47.000,1.0445) (47.500,0.8496) (48.000,0.6548) (48.500,0.4602) (49.000,0.2657) (49.500,0.0714)};
\addplot[phPsyH] coordinates {(20.706,15.4349) (21.000,15.3149) (21.500,15.1110) (22.000,14.9072) (22.500,14.7036) (23.000,14.5002) (23.500,14.2969) (24.000,14.0937) (24.500,13.8907) (25.000,13.6879) (25.500,13.4852) (26.000,13.2827) (26.500,13.0803) (27.000,12.8780) (27.500,12.6760) (28.000,12.4740) (28.500,12.2722) (29.000,12.0706) (29.500,11.8691) (30.000,11.6678) (30.500,11.4666) (31.000,11.2655) (31.500,11.0646) (32.000,10.8639) (32.500,10.6633) (33.000,10.4628) (33.500,10.2625) (34.000,10.0624) (34.500,9.8623) (35.000,9.6625) (35.500,9.4627) (36.000,9.2631) (36.500,9.0637) (37.000,8.8644) (37.500,8.6653) (38.000,8.4663) (38.500,8.2674) (39.000,8.0687) (39.500,7.8701) (40.000,7.6717) (40.500,7.4734) (41.000,7.2752) (41.500,7.0772) (42.000,6.8793) (42.500,6.6816) (43.000,6.4840) (43.500,6.2866) (44.000,6.0893) (44.500,5.8921) (45.000,5.6951) (45.500,5.4982) (46.000,5.3015) (46.500,5.1049) (47.000,4.9084) (47.500,4.7121) (48.000,4.5159) (48.500,4.3199) (49.000,4.1240) (49.500,3.9282) (50.000,3.7326)};
\addplot[phPsyH] coordinates {(23.399,18.2719) (23.500,18.2307) (24.000,18.0261) (24.500,17.8216) (25.000,17.6172) (25.500,17.4130) (26.000,17.2089) (26.500,17.0050) (27.000,16.8013) (27.500,16.5977) (28.000,16.3942) (28.500,16.1909) (29.000,15.9878) (29.500,15.7848) (30.000,15.5820) (30.500,15.3793) (31.000,15.1768) (31.500,14.9744) (32.000,14.7721) (32.500,14.5700) (33.000,14.3681) (33.500,14.1663) (34.000,13.9647) (34.500,13.7632) (35.000,13.5618) (35.500,13.3606) (36.000,13.1595) (36.500,12.9586) (37.000,12.7579) (37.500,12.5573) (38.000,12.3568) (38.500,12.1565) (39.000,11.9563) (39.500,11.7562) (40.000,11.5563) (40.500,11.3566) (41.000,11.1570) (41.500,10.9575) (42.000,10.7582) (42.500,10.5590) (43.000,10.3600) (43.500,10.1611) (44.000,9.9624) (44.500,9.7638) (45.000,9.5653) (45.500,9.3670) (46.000,9.1688) (46.500,8.9708) (47.000,8.7729) (47.500,8.5751) (48.000,8.3775) (48.500,8.1800) (49.000,7.9827) (49.500,7.7855) (50.000,7.5884)};
\addplot[phPsyH] coordinates {(25.817,21.2113) (26.000,21.1359) (26.500,20.9305) (27.000,20.7253) (27.500,20.5201) (28.000,20.3152) (28.500,20.1104) (29.000,19.9057) (29.500,19.7012) (30.000,19.4969) (30.500,19.2927) (31.000,19.0887) (31.500,18.8848) (32.000,18.6810) (32.500,18.4774) (33.000,18.2740) (33.500,18.0707) (34.000,17.8676) (34.500,17.6646) (35.000,17.4618) (35.500,17.2591) (36.000,17.0566) (36.500,16.8542) (37.000,16.6519) (37.500,16.4498) (38.000,16.2479) (38.500,16.0461) (39.000,15.8445) (39.500,15.6430) (40.000,15.4416) (40.500,15.2404) (41.000,15.0393) (41.500,14.8384) (42.000,14.6376) (42.500,14.4370) (43.000,14.2365) (43.500,14.0362) (44.000,13.8360) (44.500,13.6359) (45.000,13.4360) (45.500,13.2363) (46.000,13.0366) (46.500,12.8372) (47.000,12.6378) (47.500,12.4386) (48.000,12.2396) (48.500,12.0407) (49.000,11.8419) (49.500,11.6433) (50.000,11.4448)};
\addplot[phPsyH] coordinates {(28.004,24.2352) (28.500,24.0305) (29.000,23.8244) (29.500,23.6183) (30.000,23.4125) (30.500,23.2068) (31.000,23.0012) (31.500,22.7958) (32.000,22.5906) (32.500,22.3855) (33.000,22.1806) (33.500,21.9758) (34.000,21.7712) (34.500,21.5667) (35.000,21.3624) (35.500,21.1582) (36.000,20.9542) (36.500,20.7503) (37.000,20.5466) (37.500,20.3430) (38.000,20.1396) (38.500,19.9363) (39.000,19.7332) (39.500,19.5302) (40.000,19.3274) (40.500,19.1247) (41.000,18.9222) (41.500,18.7198) (42.000,18.5176) (42.500,18.3155) (43.000,18.1135) (43.500,17.9117) (44.000,17.7101) (44.500,17.5086) (45.000,17.3072) (45.500,17.1060) (46.000,16.9049) (46.500,16.7040) (47.000,16.5032) (47.500,16.3026) (48.000,16.1021) (48.500,15.9018) (49.000,15.7016) (49.500,15.5015) (50.000,15.3016)};
\addplot[phPsyH] coordinates {(29.998,27.3296) (30.000,27.3288) (30.500,27.1215) (31.000,26.9145) (31.500,26.7076) (32.000,26.5008) (32.500,26.2942) (33.000,26.0878) (33.500,25.8815) (34.000,25.6754) (34.500,25.4694) (35.000,25.2636) (35.500,25.0579) (36.000,24.8524) (36.500,24.6470) (37.000,24.4418) (37.500,24.2368) (38.000,24.0319) (38.500,23.8271) (39.000,23.6225) (39.500,23.4181) (40.000,23.2138) (40.500,23.0096) (41.000,22.8056) (41.500,22.6017) (42.000,22.3980) (42.500,22.1945) (43.000,21.9911) (43.500,21.7878) (44.000,21.5847) (44.500,21.3818) (45.000,21.1789) (45.500,20.9763) (46.000,20.7737) (46.500,20.5714) (47.000,20.3691) (47.500,20.1671) (48.000,19.9651) (48.500,19.7633) (49.000,19.5617) (49.500,19.3602) (50.000,19.1588)};
\addplot[phPsyH] coordinates {(33.000,29.9956) (33.500,29.7878) (34.000,29.5802) (34.500,29.3727) (35.000,29.1654) (35.500,28.9582) (36.000,28.7512) (36.500,28.5443) (37.000,28.3376) (37.500,28.1311) (38.000,27.9247) (38.500,27.7185) (39.000,27.5124) (39.500,27.3064) (40.000,27.1006) (40.500,26.8950) (41.000,26.6895) (41.500,26.4842) (42.000,26.2790) (42.500,26.0740) (43.000,25.8691) (43.500,25.6644) (44.000,25.4598) (44.500,25.2554) (45.000,25.0511) (45.500,24.8470) (46.000,24.6430) (46.500,24.4392) (47.000,24.2355) (47.500,24.0320) (48.000,23.8286) (48.500,23.6253) (49.000,23.4222) (49.500,23.2193) (50.000,23.0165)};
\addplot[phPsyH] coordinates {(42.500,29.9540) (43.000,29.7477) (43.500,29.5415) (44.000,29.3354) (44.500,29.1296) (45.000,28.9238) (45.500,28.7182) (46.000,28.5128) (46.500,28.3075) (47.000,28.1023) (47.500,27.8974) (48.000,27.6925) (48.500,27.4878) (49.000,27.2833) (49.500,27.0789) (50.000,26.8746)};
\addplot[phPsyWB] coordinates {(-5.000,2.4863) (-4.000,2.1301) (-3.000,1.7743) (-2.000,1.4191) (-1.000,1.0643) (0.000,0.7100) (1.000,0.3561) (2.000,0.0027)};
\addplot[phPsyWB] coordinates {(0.000,3.7900) (1.000,3.3920) (2.000,3.0202) (3.000,2.6613) (4.000,2.3149) (5.000,1.9816) (6.000,1.6535) (7.000,1.2991) (8.000,0.9421) (9.000,0.5816) (10.000,0.2197)};
\addplot[phPsyWB] coordinates {(5.000,5.4247) (6.000,5.0165) (7.000,4.6090) (8.000,4.2022) (9.000,3.7959) (10.000,3.3903) (11.000,2.9852) (12.000,2.5808) (13.000,2.1770) (14.000,1.7738) (15.000,1.3713) (16.000,0.9693) (17.000,0.5679) (18.000,0.1671)};
\addplot[phPsyWB] coordinates {(10.000,7.6626) (11.000,7.2508) (12.000,6.8396) (13.000,6.4290) (14.000,6.0191) (15.000,5.6097) (16.000,5.2010) (17.000,4.7929) (18.000,4.3855) (19.000,3.9786) (20.000,3.5724) (21.000,3.1668) (22.000,2.7618) (23.000,2.3573) (24.000,1.9535) (25.000,1.5503) (26.000,1.1477) (27.000,0.7457) (28.000,0.3443)};
\addplot[phPsyWB] coordinates {(15.000,10.6938) (16.000,10.2775) (17.000,9.8619) (18.000,9.4470) (19.000,9.0327) (20.000,8.6190) (21.000,8.2060) (22.000,7.7936) (23.000,7.3818) (24.000,6.9707) (25.000,6.5601) (26.000,6.1502) (27.000,5.7409) (28.000,5.3322) (29.000,4.9242) (30.000,4.5167) (31.000,4.1099) (32.000,3.7036) (33.000,3.2980) (34.000,2.8930) (35.000,2.4885) (36.000,2.0847) (37.000,1.6814) (38.000,1.2788) (39.000,0.8767) (40.000,0.4752) (41.000,0.0743)};
\addplot[phPsyWB] coordinates {(20.000,14.7605) (21.000,14.3390) (22.000,13.9182) (23.000,13.4980) (24.000,13.0785) (25.000,12.6596) (26.000,12.2414) (27.000,11.8238) (28.000,11.4069) (29.000,10.9905) (30.000,10.5749) (31.000,10.1598) (32.000,9.7454) (33.000,9.3316) (34.000,8.9185) (35.000,8.5059) (36.000,8.0940) (37.000,7.6827) (38.000,7.2720) (39.000,6.8619) (40.000,6.4525) (41.000,6.0436) (42.000,5.6354) (43.000,5.2277) (44.000,4.8207) (45.000,4.4142) (46.000,4.0083) (47.000,3.6031) (48.000,3.1984) (49.000,2.7943) (50.000,2.3909)};
\addplot[phPsyWB] coordinates {(25.000,20.1734) (26.000,19.7456) (27.000,19.3184) (28.000,18.8918) (29.000,18.4660) (30.000,18.0408) (31.000,17.6162) (32.000,17.1923) (33.000,16.7691) (34.000,16.3465) (35.000,15.9246) (36.000,15.5033) (37.000,15.0827) (38.000,14.6627) (39.000,14.2434) (40.000,13.8246) (41.000,13.4066) (42.000,12.9891) (43.000,12.5723) (44.000,12.1561) (45.000,11.7405) (46.000,11.3255) (47.000,10.9112) (48.000,10.4975) (49.000,10.0843) (50.000,9.6718)};
\addplot[phPsyWB] coordinates {(30.000,27.3329) (31.000,26.8971) (32.000,26.4621) (33.000,26.0277) (34.000,25.5940) (35.000,25.1610) (36.000,24.7287) (37.000,24.2971) (38.000,23.8661) (39.000,23.4359) (40.000,23.0062) (41.000,22.5773) (42.000,22.1490) (43.000,21.7213) (44.000,21.2944) (45.000,20.8680) (46.000,20.4423) (47.000,20.0173) (48.000,19.5929) (49.000,19.1692) (50.000,18.7461)};
\addplot[phPsyV] coordinates {(-5.500,1.8844) (-5.000,0.7137)};
\addplot[phPsyV] coordinates {(1.000,3.0696) (1.500,1.9245) (2.000,0.7839)};
\addplot[phPsyV] coordinates {(7.000,5.3267) (7.500,4.2018) (8.000,3.0814) (8.500,1.9652) (9.000,0.8534)};
\addplot[phPsyV] coordinates {(12.500,8.6000) (13.000,7.4907) (13.500,6.3856) (14.000,5.2846) (14.500,4.1878) (15.000,3.0951) (15.500,2.0065) (16.000,0.9219)};
\addplot[phPsyV] coordinates {(18.000,11.7509) (18.500,10.6569) (19.000,9.5669) (19.500,8.4810) (20.000,7.3990) (20.500,6.3210) (21.000,5.2469) (21.500,4.1768) (22.000,3.1105) (22.500,2.0480) (23.000,0.9893)};
\addplot[phPsyV] coordinates {(22.500,16.9554) (23.000,15.8686) (23.500,14.7857) (24.000,13.7068) (24.500,12.6318) (25.000,11.5606) (25.500,10.4932) (26.000,9.4297) (26.500,8.3700) (27.000,7.3140) (27.500,6.2617) (28.000,5.2132) (28.500,4.1683) (29.000,3.1271) (29.500,2.0895) (30.000,1.0556) (30.500,0.0252)};
\addplot[phPsyV] coordinates {(27.000,22.0052) (27.500,20.9258) (28.000,19.8503) (28.500,18.7785) (29.000,17.7106) (29.500,16.6465) (30.000,15.5860) (30.500,14.5294) (31.000,13.4764) (31.500,12.4271) (32.000,11.3814) (32.500,10.3394) (33.000,9.3009) (33.500,8.2661) (34.000,7.2348) (34.500,6.2071) (35.000,5.1828) (35.500,4.1621) (36.000,3.1448) (36.500,2.1310) (37.000,1.1206) (37.500,0.1136)};
\addplot[phPsyV] coordinates {(31.000,27.9821) (31.500,26.9065) (32.000,25.8346) (32.500,24.7665) (33.000,23.7021) (33.500,22.6415) (34.000,21.5845) (34.500,20.5312) (35.000,19.4815) (35.500,18.4355) (36.000,17.3930) (36.500,16.3541) (37.000,15.3187) (37.500,14.2869) (38.000,13.2586) (38.500,12.2337) (39.000,11.2124) (39.500,10.1944) (40.000,9.1799) (40.500,8.1688) (41.000,7.1610) (41.500,6.1566) (42.000,5.1556) (42.500,4.1578) (43.000,3.1634) (43.500,2.1722) (44.000,1.1843) (44.500,0.1997)};
\addplot[phPsyV] coordinates {(37.000,29.5403) (37.500,28.4834) (38.000,27.4302) (38.500,26.3805) (39.000,25.3344) (39.500,24.2919) (40.000,23.2529) (40.500,22.2174) (41.000,21.1854) (41.500,20.1568) (42.000,19.1317) (42.500,18.1101) (43.000,17.0918) (43.500,16.0769) (44.000,15.0654) (44.500,14.0572) (45.000,13.0524) (45.500,12.0508) (46.000,11.0526) (46.500,10.0576) (47.000,9.0659) (47.500,8.0774) (48.000,7.0921) (48.500,6.1100) (49.000,5.1310) (49.500,4.1553) (50.000,3.1827)};
\addplot[phPsyPt] coordinates {(20.00,7.2937)};
\addplot[phPsyPt] coordinates {(25.00,11.9490)};
\addplot[phPsyPt] coordinates {(30.00,18.8840)};
\addplot[phPsyPt] coordinates {(15.00,10.6938)};
\node[phPsyRHlab, fill=white, anchor=south east] at (axis cs:49.70,7.851) {10\,\%};
\node[phPsyRHlab, fill=white, anchor=south east] at (axis cs:49.70,15.639) {20\,\%};
\node[phPsyRHlab, fill=white, anchor=south east] at (axis cs:49.70,23.623) {30\,\%};
\node[phPsyRHlab, anchor=south] at (axis cs:48.73,30.150) {40\,\%};
\node[phPsyRHlab, anchor=south] at (axis cs:44.35,30.150) {50\,\%};
\node[phPsyRHlab, anchor=south] at (axis cs:40.87,30.150) {60\,\%};
\node[phPsyRHlab, anchor=south] at (axis cs:37.99,30.150) {70\,\%};
\node[phPsyRHlab, anchor=south] at (axis cs:35.55,30.150) {80\,\%};
\node[phPsyRHlab, anchor=south] at (axis cs:33.43,30.150) {90\,\%};
\node[phPsyHlab, anchor=south east] at (axis cs:-5.96,2.478) {0};
\node[phPsyHlab, anchor=south east] at (axis cs:0.11,4.026) {10};
\node[phPsyHlab, anchor=south east] at (axis cs:5.48,5.841) {20};
\node[phPsyHlab, anchor=south east] at (axis cs:10.08,7.957) {30};
\node[phPsyHlab, anchor=south east] at (axis cs:14.03,10.318) {40};
\node[phPsyHlab, anchor=south east] at (axis cs:17.48,12.874) {50};
\node[phPsyHlab, anchor=south east] at (axis cs:20.51,15.585) {60};
\node[phPsyHlab, anchor=south east] at (axis cs:23.20,18.422) {70};
\node[phPsyHlab, anchor=south east] at (axis cs:25.62,21.361) {80};
\node[phPsyHlab, anchor=south east] at (axis cs:27.80,24.385) {90};
\node[phPsyHlab, anchor=south east] at (axis cs:29.80,27.480) {100};
\node[phPsyHlab, fill=white, inner sep=0.8pt] at (axis cs:36.00,28.751) {110};
\node[phPsyHlab, fill=white, inner sep=0.8pt] at (axis cs:45.50,28.718) {120};
\node[phPsyWBlab, fill=white] at (axis cs:-1.50,1.242) {$-$5\,℃};
\node[phPsyWBlab, fill=white] at (axis cs:7.50,1.122) {0\,℃};
\node[phPsyWBlab, fill=white] at (axis cs:15.50,1.170) {5\,℃};
\node[phPsyWBlab, fill=white] at (axis cs:25.50,1.349) {10\,℃};
\node[phPsyWBlab, fill=white] at (axis cs:38.50,1.078) {15\,℃};
\node[phPsyWBlab, fill=white] at (axis cs:47.50,3.401) {20\,℃};
\node[phPsyWBlab, fill=white] at (axis cs:47.50,10.704) {25\,℃};
\node[phPsyWBlab, fill=white] at (axis cs:47.50,19.805) {30\,℃};
\node[phPsyVlab, anchor=north, yshift=-4.6mm] at (axis cs:-5.00,0) {0.76};
\node[phPsyVlab, anchor=north, yshift=-4.6mm] at (axis cs:2.00,0) {0.78};
\node[phPsyVlab, anchor=north, yshift=-4.6mm] at (axis cs:9.00,0) {0.80};
\node[phPsyVlab, anchor=north, yshift=-4.6mm] at (axis cs:16.00,0) {0.82};
\node[phPsyVlab, anchor=north, yshift=-4.6mm] at (axis cs:23.00,0) {0.84};
\node[phPsyVlab, anchor=north, yshift=-4.6mm] at (axis cs:30.50,0) {0.86};
\node[phPsyVlab, anchor=north, yshift=-4.6mm] at (axis cs:37.50,0) {0.88};
\node[phPsyVlab, anchor=north, yshift=-4.6mm] at (axis cs:44.50,0) {0.90};
\node[phPsyVlab, anchor=north, yshift=-4.6mm] at (axis cs:50.00,0) {0.92};
\node[phPsyVlab, anchor=north east, yshift=-4.6mm] at (axis cs:-10.30,0) {比体積〔m$^3$/kg(DA)〕};
\node[phPsyPtlab, anchor=north west] at (axis cs:20.30,7.194) {A};
\node[phPsyPtlab, anchor=east] at (axis cs:24.70,11.899) {B};
\node[phPsyPtlab, anchor=east] at (axis cs:29.70,18.834) {C};
\node[phPsyPtlab, anchor=north west] at (axis cs:15.30,10.594) {D};
\end{axis}
% 修2-6（Q6-14）：右のふちにも絶対湿度の目もりと名前（見開きの右のページで、左のページまで線をたどらずに x を読む）
\begin{axis}[at={(psy.south west)}, anchor=south west, width=#1, height=#2, scale only axis, xmin=-10, xmax=50, ymin=0, ymax=30,
  hide x axis, axis y line*=right, axis line style={draw=none}, ytick={0,2,...,30}, yticklabel pos=right, ylabel near ticks,
  tick label style={font=\sffamily\footnotesize, text=bkInk}, major tick length=0pt,
  ylabel={絶対湿度 $x$〔g/kg(DA)〕}, label style={font=\sffamily\small, text=bkInk}, scaled ticks=false]
\addplot[draw=none] coordinates {(-10,0)};
\end{axis}
% 図の左はしから軸の左のふちまでの長さを \ph@psyoff に残す（見開きの切れ目を、軸の上の決まった温度に置くため）
\pgfextractx{\pgf@xa}{\pgfpointanchor{psy}{south west}}\pgfextractx{\pgf@xb}{\pgfpointanchor{current bounding box}{south west}}%
\global\ph@psyoff=\dimexpr\pgf@xa-\pgf@xb\relax
\end{tikzpicture}}

\newsavebox\ph@psybox
\newlength\ph@psyw
\newlength\ph@psyoff
\newlength\ph@psysplit
\newcommand\phPsychroTop[1]{\noindent\begin{minipage}[t][11mm][t]{\linewidth}#1\end{minipage}\par}
%% 修2-6（Q6-13）：空気線図と湿り空気の表は、ページの偶奇で組む順が入れかわるので、番号は組む順に 4・5 とつける（#1）
\newcommand\phPsychroSpread[1]{%
  % 修2-6（Q6-14）：右のページの右のふちの外に、絶対湿度の目もりと名前（約 11 mm）を置くので、図のはばを
  % 「17.5 ℃から右（65/120）が、本文のはば −11 mm に入る」長さにする（前は 2\linewidth−32mm で、右に 6 mm しかなかった）
  \setlength\ph@psyw{\dimexpr(\linewidth-11mm)*120/65\relax}%
  \sbox\ph@psybox{\phPsychroAxis{\ph@psyw}{130mm}}%
  % 見開きの切れ目は、乾球温度 17.5 ℃の所（目もりの 15 と 20 の間。本文の点 A〜D をさける）。軸は −10〜50 ℃
  \setlength\ph@psysplit{\dimexpr\ph@psyoff+\ph@psyw*55/120\relax}%
  \clearpage
  \phPsychroTop{\phbkhead{#1　空気線図（湿り空気線図、1気圧）}}
  \noindent\begin{tikzpicture}\path[use as bounding box] (0,0) rectangle (\linewidth,\ht\ph@psybox+\dp\ph@psybox);
    \clip (-2mm,-2mm) rectangle (\linewidth,\ht\ph@psybox+\dp\ph@psybox+2mm);
    \node[anchor=south west, inner sep=0pt] at (\linewidth-\ph@psysplit,0) {\usebox\ph@psybox};\end{tikzpicture}\par
  \vspace{7mm}{\footnotesize\color{bkGray}見開きの図です（右のページへつづく）。横は乾球温度 $\theta$〔℃〕、縦は絶対湿度 $x$〔g/kg(DA)〕（目もりは、このページの左のふちと、右のページの右のふち）。相対湿度の数は線の右はし（上のふちの外か、右のふちのすぐ内側）、比体積の数は下のふちの外。斜めの赤い線が比エンタルピー、ほぼたての点線が比体積。CoolProp 8.0 の湿り空気の式（HAPropsSI、全圧\qty{101.325}{kPa}）で計算した。本文の式 $x=0.622\,p_{\mathrm v}/(p-p_{\mathrm v})$ は理想気体の式で、$x$ が 0.4\,\% ほど小さく出る（熱-9）。\par}
  \clearpage
  \phPsychroTop{\hfill{\sffamily\footnotesize\color{bkGray}空気線図（左のページからつづく）}}
  \noindent\begin{tikzpicture}\path[use as bounding box] (0,0) rectangle (\linewidth,\ht\ph@psybox+\dp\ph@psybox);
    \clip (0,-2mm) rectangle (\linewidth+16mm,\ht\ph@psybox+\dp\ph@psybox+2mm);
    \node[anchor=south west, inner sep=0pt] at (-\ph@psysplit,0) {\usebox\ph@psybox};\end{tikzpicture}\par
  \vspace{7mm}
  {\footnotesize
  \tikz[baseline=-0.5ex]\draw[bkBlue!80!black, line width=1.6pt] (0,0) -- (5mm,0);\ 飽和（相対湿度100\,\%）\quad
  \tikz[baseline=-0.5ex]\draw[bkBlue!80!black, line width=0.6pt] (0,0) -- (5mm,0);\ 相対湿度 $\varphi$（10\,\% ごと）\quad
  \tikz[baseline=-0.5ex]\draw[bkRed!75!black, line width=0.45pt] (0,0) -- (5mm,0);\ 比エンタルピー $h$〔kJ/kg(DA)〕（10 ごと。数は飽和線の左上。110・120 は線の上）\par
  \tikz[baseline=-0.5ex]\draw[bkGreen!70!black, line width=0.45pt, dash pattern=on 2.4pt off 1.6pt] (0,0) -- (5mm,0);\ 湿球温度（\qty{5}{\celsius}ごと。飽和線の上の温度から右下へ。数は「20\,℃」のように ℃ をつけて、線の下のはしの近く。比エンタルピーの線とほぼ平行）\quad
  \tikz[baseline=-0.5ex]\draw[bkPurple!75!black, line width=0.45pt, dash pattern=on 0.8pt off 1.2pt] (0,0) -- (5mm,0);\ 比体積〔m$^3$/kg(DA)〕\par
  点 A：\qty{20}{\celsius}・50\,\%（露点 \qty{9.3}{\celsius}、熱-9・熱-11 の窓の結露）、B：\qty{25}{\celsius}・60\,\%（湿球 \qty{19.5}{\celsius}、熱-9 例題10）、C→D：\qty{30}{\celsius}・70\,\% を\qty{15}{\celsius}の飽和まで冷やす（熱-9 例題4）。\par
  \par}
  \clearpage}
\newcommand\phMoistPage[1]{%
  \phbkhead{#1　湿り空気の表（空気線図の抜き出し、1気圧）}
  {\small 空気線図の点を、数で読んだ表です。熱-9 の例題の点をふくみます。絶対湿度 $x$ と比エンタルピー $h$ は、乾いた空気\qty{1}{kg}あたりの値です。露点は、その空気を冷やしたとき、水滴がつき始める温度です（熱-9・熱-11 の結露）。\par}\vspace{1mm}
  %% 自動生成（tools/mksteam2.py）。直接編集しないこと。値は IAPWS-IF97 の式（iapws 1.5.5）・CoolProp 8.0.0
\begin{longtblr}{colspec={Q[c,m] Q[c,m] Q[r,m] Q[r,m] Q[r,m] Q[r,m] Q[r,m]}, width=\linewidth, rowhead=1, colsep=3.2pt, rowsep=0.35pt, hlines={0.3pt,bkLine}, vline{2}={0.4pt,bkLine}, row{1-1}={bg=bkVol!14, font=\sffamily\bfseries\scriptsize, halign=c}, cells={font=\footnotesize}}
{温度\\$\theta$〔℃〕} & {相対湿度\\$\varphi$〔\%〕} & {絶対湿度 $x$\\〔kg/kg(DA)〕} & {比エンタルピー $h$\\〔kJ/kg(DA)〕} & {比体積\\〔m$^3$/kg(DA)〕} & {湿球温度\\〔℃〕} & {露点\\〔℃〕} \\
0 & 80 & 0.0030 & 7.6 & 0.777 & $-$1.2 & $-$2.7 \\
5 & 60 & 0.0032 & 13.2 & 0.792 & 2.1 & $-$1.9 \\
10 & 60 & 0.0046 & 21.6 & 0.808 & 6.5 & 2.6 \\
15 & 100 & 0.0107 & 42.1 & 0.830 & 15.0 & 15.0 \\
20 & 21 & 0.0030 & 27.8 & 0.834 & 9.4 & $-$2.6 \\
20 & 40 & 0.0058 & 34.9 & 0.838 & 12.3 & 6.0 \\
20 & 50 & 0.0073 & 38.6 & 0.840 & 13.8 & 9.3 \\
20 & 60 & 0.0088 & 42.4 & 0.842 & 15.1 & 12.0 \\
22 & 40 & 0.0066 & 38.9 & 0.845 & 13.9 & 7.8 \\
25 & 50 & 0.0099 & 50.4 & 0.858 & 17.9 & 13.9 \\
25 & 60 & 0.0119 & 55.6 & 0.861 & 19.5 & 16.7 \\
26 & 50 & 0.0105 & 53.0 & 0.862 & 18.7 & 14.8 \\
28 & 60 & 0.0143 & 64.7 & 0.872 & 22.1 & 19.5 \\
28 & 70 & 0.0168 & 70.9 & 0.876 & 23.7 & 22.0 \\
30 & 50 & 0.0134 & 64.4 & 0.877 & 22.0 & 18.5 \\
30 & 70 & 0.0189 & 78.4 & 0.885 & 25.5 & 23.9 \\
32 & 60 & 0.0181 & 78.6 & 0.889 & 25.6 & 23.3 \\
35 & 40 & 0.0142 & 71.6 & 0.893 & 23.9 & 19.4 \\
35 & 60 & 0.0215 & 90.5 & 0.903 & 28.2 & 26.1 \\
\end{longtblr}

  {\footnotesize\color{bkGray}CoolProp 8.0（HAPropsSI、\qty{101.325}{kPa}）で計算。中学の飽和水蒸気量（\qty{25}{\celsius}で約\qty{23}{g/m^3}）は、IAPWS-IF97 の飽和蒸気圧と理想気体の式から求めた値（熱-9）。\par}}
%% 空気線図の見開きは、偶数（左）のページから始める。湿り空気の表（約半ページ）は、過熱蒸気表のあとに続けると、次のページの上で終わる。
%%   過熱蒸気表が偶数のページで終わったとき：湿り空気の表を続けて組むと、次の偶数のページから見開きになる。
%%   奇数のページで終わったとき：改ページして見開きを先に組み、湿り空気の表はそのあと。
%% それでも奇数のページから始まるときだけ、メモのページが1枚入る（BKNOTE がログに出る。並びを見直す）。
\ifodd\value{page}
  \clearpage
  \ifodd\value{page}\typeout{BKNOTE 空気線図の前にメモのページが入った}\bkmemopage\fi
  \phPsychroSpread{4}
  \phMoistPage{5}
\else
  \phMoistPage{4}
  \clearpage
  \ifodd\value{page}\typeout{BKNOTE 空気線図の前にメモのページが入った}\bkmemopage\fi
  \phPsychroSpread{5}
\fi

\phbkhead{6　冷媒の飽和表}
{\small 家庭用エアコンの R32、カーエアコン・業務用の R134a、家庭用冷蔵庫の R600a（イソブタン）です（熱-9）。CoolProp 8.0 の各冷媒の状態式で計算しました。\par}\vspace{1mm}
%% 自動生成（tools/mksteam2.py）。直接編集しないこと。値は IAPWS-IF97 の式（iapws 1.5.5）・CoolProp 8.0.0
\phbkhead{R32（ジフルオロメタン。家庭用エアコン）の飽和表}
{\footnotesize\color{bkGray}臨界点 \qty{78.1}{\celsius}・\qty{5.78}{MPa}。\par}\vspace{0.5mm}
\begin{longtblr}{colspec={Q[c,m,wd=9mm] Q[r,m] Q[r,m] Q[r,m] Q[r,m] Q[r,m] Q[r,m] Q[r,m]}, width=\linewidth, rowhead=1, colsep=3.2pt, rowsep=0.35pt, hlines={0.3pt,bkLine}, vline{2}={0.4pt,bkLine}, row{1-1}={bg=bkVol!14, font=\sffamily\bfseries\scriptsize, halign=c}, cells={font=\scriptsize}}
{温度\\$\theta$〔℃〕} & {圧力\\$p$〔MPa〕} & {液の密度\\$\rho'$〔kg/m$^3$〕} & {$v''$\\〔m$^3$/kg〕} & {$h'$\\〔kJ/kg〕} & {$h''$\\〔kJ/kg〕} & {$s'$\\〔kJ/(kg·K)〕} & {$s''$\\〔kJ/(kg·K)〕} \\
$-$40 & 0.1774 & 1180.2 & 0.1974 & 133.2 & 502.0 & 0.7382 & 2.3200 \\
$-$35 & 0.2214 & 1165.7 & 0.1601 & 141.3 & 504.2 & 0.7723 & 2.2962 \\
$-$30 & 0.2734 & 1151.0 & 0.1309 & 149.4 & 506.3 & 0.8060 & 2.2735 \\
$-$25 & 0.3346 & 1135.9 & 0.1079 & 157.7 & 508.2 & 0.8392 & 2.2518 \\
$-$20 & 0.4058 & 1120.6 & 0.08963 & 165.9 & 510.0 & 0.8720 & 2.2310 \\
$-$15 & 0.4881 & 1104.9 & 0.07493 & 174.3 & 511.6 & 0.9044 & 2.2109 \\
$-$10 & 0.5826 & 1088.8 & 0.06301 & 182.8 & 513.0 & 0.9365 & 2.1915 \\
$-$5 & 0.6906 & 1072.2 & 0.05328 & 191.3 & 514.3 & 0.9684 & 2.1727 \\
0 & 0.8131 & 1055.3 & 0.04527 & 200.0 & 515.3 & 1.0000 & 2.1543 \\
5 & 0.9514 & 1037.7 & 0.03862 & 208.8 & 516.1 & 1.0314 & 2.1363 \\
10 & 1.107 & 1019.7 & 0.03308 & 217.7 & 516.7 & 1.0628 & 2.1185 \\
15 & 1.281 & 1000.9 & 0.02842 & 226.8 & 516.9 & 1.0940 & 2.1008 \\
20 & 1.475 & 981.4 & 0.02448 & 236.1 & 516.9 & 1.1253 & 2.0831 \\
25 & 1.690 & 961.0 & 0.02112 & 245.6 & 516.5 & 1.1566 & 2.0652 \\
30 & 1.928 & 939.6 & 0.01826 & 255.3 & 515.7 & 1.1881 & 2.0471 \\
35 & 2.190 & 917.0 & 0.01579 & 265.3 & 514.5 & 1.2198 & 2.0285 \\
40 & 2.478 & 893.0 & 0.01365 & 275.6 & 512.7 & 1.2520 & 2.0091 \\
45 & 2.795 & 867.3 & 0.01178 & 286.3 & 510.3 & 1.2847 & 1.9888 \\
50 & 3.141 & 839.3 & 0.01015 & 297.5 & 507.1 & 1.3183 & 1.9670 \\
55 & 3.520 & 808.3 & 0.008697 & 309.3 & 502.9 & 1.3531 & 1.9432 \\
60 & 3.933 & 773.3 & 0.007396 & 321.9 & 497.4 & 1.3898 & 1.9166 \\
\end{longtblr}
\phbkhead{R134a（テトラフルオロエタン。カーエアコン・冷蔵庫）の飽和表}
{\footnotesize\color{bkGray}臨界点 \qty{101.1}{\celsius}・\qty{4.06}{MPa}。\par}\vspace{0.5mm}
\begin{longtblr}{colspec={Q[c,m,wd=9mm] Q[r,m] Q[r,m] Q[r,m] Q[r,m] Q[r,m] Q[r,m] Q[r,m]}, width=\linewidth, rowhead=1, colsep=3.2pt, rowsep=0.35pt, hlines={0.3pt,bkLine}, vline{2}={0.4pt,bkLine}, row{1-1}={bg=bkVol!14, font=\sffamily\bfseries\scriptsize, halign=c}, cells={font=\scriptsize}}
{温度\\$\theta$〔℃〕} & {圧力\\$p$〔MPa〕} & {液の密度\\$\rho'$〔kg/m$^3$〕} & {$v''$\\〔m$^3$/kg〕} & {$h'$\\〔kJ/kg〕} & {$h''$\\〔kJ/kg〕} & {$s'$\\〔kJ/(kg·K)〕} & {$s''$\\〔kJ/(kg·K)〕} \\
$-$40 & 0.05121 & 1417.7 & 0.3611 & 148.1 & 374.0 & 0.7956 & 1.7643 \\
$-$35 & 0.06614 & 1403.1 & 0.2840 & 154.4 & 377.2 & 0.8223 & 1.7575 \\
$-$30 & 0.08438 & 1388.4 & 0.2259 & 160.8 & 380.3 & 0.8486 & 1.7515 \\
$-$25 & 0.1064 & 1373.4 & 0.1816 & 167.2 & 383.4 & 0.8746 & 1.7461 \\
$-$20 & 0.1327 & 1358.3 & 0.1474 & 173.6 & 386.6 & 0.9002 & 1.7413 \\
$-$15 & 0.1639 & 1342.8 & 0.1207 & 180.1 & 389.6 & 0.9256 & 1.7371 \\
$-$10 & 0.2006 & 1327.1 & 0.09959 & 186.7 & 392.7 & 0.9506 & 1.7334 \\
$-$5 & 0.2433 & 1311.1 & 0.08280 & 193.3 & 395.7 & 0.9754 & 1.7300 \\
0 & 0.2928 & 1294.8 & 0.06931 & 200.0 & 398.6 & 1.0000 & 1.7271 \\
5 & 0.3497 & 1278.1 & 0.05837 & 206.8 & 401.5 & 1.0243 & 1.7245 \\
10 & 0.4146 & 1261.0 & 0.04944 & 213.6 & 404.3 & 1.0485 & 1.7221 \\
15 & 0.4884 & 1243.4 & 0.04209 & 220.5 & 407.1 & 1.0724 & 1.7200 \\
20 & 0.5717 & 1225.3 & 0.03600 & 227.5 & 409.7 & 1.0962 & 1.7180 \\
25 & 0.6654 & 1206.7 & 0.03091 & 234.5 & 412.3 & 1.1199 & 1.7162 \\
30 & 0.7702 & 1187.5 & 0.02664 & 241.7 & 414.8 & 1.1435 & 1.7145 \\
35 & 0.8870 & 1167.5 & 0.02303 & 249.0 & 417.2 & 1.1670 & 1.7128 \\
40 & 1.017 & 1146.7 & 0.01997 & 256.4 & 419.4 & 1.1905 & 1.7111 \\
45 & 1.160 & 1125.1 & 0.01734 & 263.9 & 421.5 & 1.2139 & 1.7092 \\
50 & 1.318 & 1102.3 & 0.01509 & 271.6 & 423.4 & 1.2375 & 1.7072 \\
55 & 1.492 & 1078.3 & 0.01314 & 279.5 & 425.2 & 1.2611 & 1.7050 \\
60 & 1.682 & 1052.9 & 0.01144 & 287.5 & 426.6 & 1.2848 & 1.7024 \\
\end{longtblr}
\phbkhead{R600a（イソブタン。家庭用冷蔵庫）の飽和表}
{\footnotesize\color{bkGray}臨界点 \qty{134.7}{\celsius}・\qty{3.63}{MPa}。\par}\vspace{0.5mm}
\begin{longtblr}{colspec={Q[c,m,wd=9mm] Q[r,m] Q[r,m] Q[r,m] Q[r,m] Q[r,m] Q[r,m] Q[r,m]}, width=\linewidth, rowhead=1, colsep=3.2pt, rowsep=0.35pt, hlines={0.3pt,bkLine}, vline{2}={0.4pt,bkLine}, row{1-1}={bg=bkVol!14, font=\sffamily\bfseries\scriptsize, halign=c}, cells={font=\scriptsize}}
{温度\\$\theta$〔℃〕} & {圧力\\$p$〔MPa〕} & {液の密度\\$\rho'$〔kg/m$^3$〕} & {$v''$\\〔m$^3$/kg〕} & {$h'$\\〔kJ/kg〕} & {$h''$\\〔kJ/kg〕} & {$s'$\\〔kJ/(kg·K)〕} & {$s''$\\〔kJ/(kg·K)〕} \\
$-$40 & 0.02870 & 624.1 & 1.143 & 112.5 & 501.3 & 0.6549 & 2.3227 \\
$-$35 & 0.03680 & 618.9 & 0.9074 & 123.0 & 507.9 & 0.6995 & 2.3154 \\
$-$30 & 0.04662 & 613.6 & 0.7284 & 133.7 & 514.4 & 0.7437 & 2.3095 \\
$-$25 & 0.05843 & 608.3 & 0.5906 & 144.4 & 521.0 & 0.7874 & 2.3048 \\
$-$20 & 0.07248 & 602.9 & 0.4834 & 155.3 & 527.6 & 0.8306 & 2.3013 \\
$-$15 & 0.08905 & 597.4 & 0.3990 & 166.3 & 534.3 & 0.8735 & 2.2989 \\
$-$10 & 0.1085 & 591.9 & 0.3320 & 177.4 & 540.9 & 0.9160 & 2.2975 \\
$-$5 & 0.1310 & 586.3 & 0.2783 & 188.6 & 547.6 & 0.9582 & 2.2969 \\
0 & 0.1570 & 580.6 & 0.2349 & 200.0 & 554.3 & 1.0000 & 2.2972 \\
5 & 0.1867 & 574.8 & 0.1995 & 211.5 & 561.1 & 1.0415 & 2.2983 \\
10 & 0.2206 & 568.9 & 0.1704 & 223.1 & 567.8 & 1.0828 & 2.3000 \\
15 & 0.2590 & 562.9 & 0.1464 & 234.9 & 574.5 & 1.1239 & 2.3023 \\
20 & 0.3022 & 556.9 & 0.1264 & 246.9 & 581.2 & 1.1647 & 2.3051 \\
25 & 0.3507 & 550.7 & 0.1096 & 259.0 & 587.9 & 1.2053 & 2.3085 \\
30 & 0.4047 & 544.3 & 0.09542 & 271.2 & 594.6 & 1.2458 & 2.3123 \\
35 & 0.4648 & 537.8 & 0.08341 & 283.7 & 601.2 & 1.2861 & 2.3165 \\
40 & 0.5312 & 531.2 & 0.07317 & 296.3 & 607.8 & 1.3263 & 2.3211 \\
45 & 0.6044 & 524.4 & 0.06440 & 309.1 & 614.3 & 1.3664 & 2.3259 \\
50 & 0.6849 & 517.4 & 0.05683 & 322.1 & 620.8 & 1.4064 & 2.3309 \\
55 & 0.7730 & 510.2 & 0.05029 & 335.3 & 627.2 & 1.4464 & 2.3361 \\
60 & 0.8692 & 502.7 & 0.04459 & 348.7 & 633.5 & 1.4863 & 2.3414 \\
\end{longtblr}

\phbkhead{冷媒のくらべ}
{\footnotesize
\begin{longtblr}{colspec={Q[l,m,wd=20mm] Q[c,m,wd=26mm] Q[c,m,wd=19mm] Q[c,m,wd=16mm] X[l,m]}, width=\linewidth, rowhead=1, \ph@tblstyle}
冷媒 & 臨界点 & 温暖化係数 GWP & 安全区分 & 出典 \\
R32 & \qty{78.1}{\celsius}・\qty{5.78}{MPa} & 675 & A2L & 臨界点：CoolProp 8.0。GWP：IPCC 第4次評価報告書（100年値）。安全区分：ISO 817・ASHRAE 34 \\
R134a & \qty{101.1}{\celsius}・\qty{4.06}{MPa} & 1430 & A1 & 同じ \\
R600a & \qty{134.7}{\celsius}・\qty{3.63}{MPa} & 約3 & A3 & GWP は IIR などの資料（AR4 の表にない）\phyk \\
R744（CO$_2$） & \qty{31.0}{\celsius}・\qty{7.38}{MPa} & 1 & A1 & 臨界点：CoolProp 8.0（Span・Wagner の式） \\
R410A（前の冷媒） & ― & 2088 & A1 & IPCC 第4次評価報告書 \\
\end{longtblr}\par}
{\footnotesize\color{bkGray}A は毒性が低い、1 は燃えない、2L は燃えにくい（微燃性）、3 は燃えやすい。家庭用エアコンの冷媒の GWP の目標 750 以下（フロン排出抑制法の指定製品）\phyk。日本のフロン排出抑制法の告示も AR4 の GWP を使う\phyk。\par}

\phbkhead{7　気体の性質}
%% 自動生成（tools/mksteam2.py）。直接編集しないこと。値は IAPWS-IF97 の式（iapws 1.5.5）・CoolProp 8.0.0
\begin{longtblr}{colspec={Q[l,m,wd=24mm] Q[r,m] Q[r,m] Q[r,m] Q[r,m] Q[r,m] Q[r,m] Q[r,m]}, width=\linewidth, rowhead=1, colsep=3.2pt, rowsep=0.35pt, hlines={0.3pt,bkLine}, vline{2}={0.4pt,bkLine}, row{1-1}={bg=bkVol!14, font=\sffamily\bfseries\scriptsize, halign=c}, cells={font=\scriptsize}}
気体 & {モル質量 $M$\\〔g/mol〕} & {$R_{\mathrm s}=R/M$\\〔J/(kg·K)〕} & {$C_V$\\〔J/(mol·K)〕} & {$C_p$\\〔J/(mol·K)〕} & {比熱比\\$\gamma$} & {音速 $a$\\〔m/s〕} & {圧縮係数\\$Z$} \\
ヘリウム He & 4.003 & 2077 & 12.5 & 20.8 & 1.67 & 1008 & 1.0005 \\
アルゴン Ar & 39.948 & 208 & 12.5 & 20.8 & 1.67 & 319 & 0.9994 \\
水素 H$_2$ & 2.016 & 4124 & 20.5 & 28.9 & 1.41 & 1305 & 1.0006 \\
窒素 N$_2$ & 28.013 & 297 & 20.8 & 29.2 & 1.40 & 349 & 0.9998 \\
酸素 O$_2$ & 31.999 & 260 & 21.1 & 29.4 & 1.40 & 326 & 0.9994 \\
空気 & 28.965 & 287 & 20.8 & 29.2 & 1.40 & 343 & 0.9997 \\
二酸化炭素 CO$_2$ & 44.010 & 189 & 29.0 & 37.5 & 1.29 & 267 & 0.9951 \\
メタン CH$_4$ & 16.043 & 518 & 27.5 & 35.9 & 1.31 & 445 & 0.9983 \\
\end{longtblr}
\par{\footnotesize\color{bkGray}$C_V$・$C_p$・$\gamma$・$Z$ は\qty{300}{K}、音速は\qty{20}{\celsius}の値（1気圧。CoolProp 8.0 で計算）。熱-3 の表、熱-4 の音速（空気 343、ヘリウム 1008）、熱-14 の等分配と同じ値。\par}\vspace{2mm}
\phbkhead{空気の物性（1気圧）}
\begin{longtblr}{colspec={Q[c,m,wd=11mm] Q[r,m] Q[r,m] Q[r,m] Q[r,m] Q[r,m] Q[r,m] Q[r,m]}, width=\linewidth, rowhead=1, colsep=3.2pt, rowsep=0.35pt, hlines={0.3pt,bkLine}, vline{2}={0.4pt,bkLine}, row{1-1}={bg=bkVol!14, font=\sffamily\bfseries\scriptsize, halign=c}, cells={font=\scriptsize}}
{温度\\$\theta$〔℃〕} & {密度 $\rho$\\〔kg/m$^3$〕} & {比熱 $c_p$\\〔J/(kg·K)〕} & {熱伝導率 $k$\\〔W/(m·K)〕} & {粘度 $\mu$\\〔$10^{-5}$\,Pa·s〕} & {動粘度 $\nu$\\〔$10^{-6}$\,m$^2$/s〕} & {プラントル数\\Pr} & {熱拡散率 $\alpha$\\〔$10^{-6}$\,m$^2$/s〕} \\
$-$20 & 1.396 & 1005.5 & 0.0228 & 1.62 & 11.6 & 0.71 & 16.3 \\
0 & 1.293 & 1005.7 & 0.0244 & 1.72 & 13.3 & 0.71 & 18.7 \\
20 & 1.205 & 1006.1 & 0.0259 & 1.82 & 15.1 & 0.71 & 21.3 \\
27 & 1.176 & 1006.4 & 0.0264 & 1.85 & 15.8 & 0.71 & 22.3 \\
40 & 1.127 & 1006.9 & 0.0274 & 1.92 & 17.0 & 0.71 & 24.1 \\
60 & 1.060 & 1008.0 & 0.0288 & 2.01 & 19.0 & 0.70 & 27.0 \\
100 & 0.946 & 1011.2 & 0.0316 & 2.19 & 23.1 & 0.70 & 33.1 \\
200 & 0.746 & 1025.0 & 0.0382 & 2.60 & 34.9 & 0.70 & 50.0 \\
\end{longtblr}
\phbkhead{水の物性（1気圧の液）}
\begin{longtblr}{colspec={Q[c,m,wd=10mm] Q[r,m] Q[r,m] Q[r,m] Q[r,m] Q[r,m] Q[r,m] Q[r,m]}, width=\linewidth, rowhead=1, colsep=3.2pt, rowsep=0.35pt, hlines={0.3pt,bkLine}, vline{2}={0.4pt,bkLine}, row{1-1}={bg=bkVol!14, font=\sffamily\bfseries\scriptsize, halign=c}, cells={font=\scriptsize}}
{温度\\$\theta$〔℃〕} & {密度 $\rho$\\〔kg/m$^3$〕} & {比熱 $c_p$\\〔J/(kg·K)〕} & {熱伝導率 $k$\\〔W/(m·K)〕} & {粘度 $\mu$\\〔mPa·s〕} & {プラントル数\\Pr} & {熱拡散率 $\alpha$\\〔$10^{-7}$\,m$^2$/s〕} & {体膨張率 $\beta$\\〔$10^{-4}$/K〕} \\
0.01 & 999.8 & 4219 & 0.556 & 1.791 & 13.60 & 1.317 & $-$0.68 \\
10 & 999.7 & 4195 & 0.579 & 1.306 & 9.47 & 1.380 & 0.88 \\
20 & 998.2 & 4184 & 0.598 & 1.002 & 7.01 & 1.432 & 2.07 \\
25 & 997.0 & 4181 & 0.607 & 0.890 & 6.14 & 1.455 & 2.57 \\
40 & 992.2 & 4179 & 0.628 & 0.653 & 4.34 & 1.516 & 3.85 \\
60 & 983.2 & 4185 & 0.651 & 0.466 & 3.00 & 1.582 & 5.23 \\
80 & 971.8 & 4197 & 0.667 & 0.354 & 2.23 & 1.635 & 6.41 \\
99 & 959.1 & 4215 & 0.677 & 0.285 & 1.77 & 1.674 & 7.45 \\
\end{longtblr}
\par{\footnotesize\color{bkGray}空気・水の値は CoolProp 8.0（空気：Lemmon ほかの状態式、水：IAPWS-95 と IAPWS の粘度・熱伝導率の式）で計算した。本文の目安：空気 $c_p=\qty{1.005}{kJ/(kg.K)}$（\qty{20}{\celsius}で 1.006。熱-3・熱-4・熱-6・熱-7・熱-8）・$R_{\mathrm s}=\qty{287}{J/(kg.K)}$（熱-3・熱-4・熱-6・熱-7）・$\gamma=1.40$・$k=0.026$・$\mu=1.8\times10^{-5}$、水 $c=\qty{4.18}{kJ/(kg.K)}$（本文は 4.2）・$k=0.60$。\par}


\phbkhead{8　液体・固体の熱の性質（室温）}
{\footnotesize
\begin{longtblr}{colspec={Q[l,m,wd=25mm] Q[c,m] Q[c,m] Q[c,m] Q[c,m] Q[c,m] Q[c,m,wd=15mm]}, width=\linewidth, rowhead=1, \ph@tblstyle, colsep=1.8pt}
物 & {密度 $\rho$\\〔kg/m$^3$〕} & {比熱 $c$\\〔kJ/(kg·K)〕} & {熱伝導率 $k$\\〔W/(m·K)〕} & {熱拡散率 $\alpha$\\〔m$^2$/s〕} & {線膨張係数\\〔1/K〕} & 出典 \\
銅 & 8900 & 0.39（0.385） & 約400 & $1.2\times10^{-4}$ & $1.7\times10^{-5}$ & ア・イ・ウ \\
アルミニウム & 2700 & 0.90（0.897） & 237 & $1.0\times10^{-4}$ & $2.3\times10^{-5}$ & ア・ウ \\
鉄・鋼 & 7850 & 0.45（0.449） & 鋼 約50 & 鉄 $2.3\times10^{-5}$ & $1.2\times10^{-5}$ & ア・イ・ウ・エ \\
ステンレス鋼 & 約7900 & 約0.50 & 16 & $4\times10^{-6}$ & ― & ウ・エ \\
ねずみ鋳鉄（ブレーキの円板） & 約7200 & 約0.46 & ― & ― & ― & カ（熱-2） \\
ガラス & 約2500 & 約0.84 & 約1.0 & ― & ― & イ \\
コンクリート & 2300 & 0.88 & 1.6 & $7.9\times10^{-7}$ & ― & ウ（熱-11） \\
れんが & ― & ― & 0.6 & ― & ― & ウ・エ \\
木 & ― & ― & 約0.15 & ― & ― & イ \\
グラスウール & 16 & 0.84 & 0.04 & ― & ― & ウ・カ \\
土（湿ったもの） & ― & ― & ― & $5\times10^{-7}$ & ― & キ（熱-11） \\
氷（\qty{0}{\celsius}） & 917 & 2.1 & ― & ― & ― & エ（熱-1） \\
水（\qty{20}{\celsius}） & 998.2 & 4.18 & 0.60 & $1.4\times10^{-7}$ & ― & オ \\
エタノール & ― & 2.4 & ― & ― & ― & ア \\
作動油（VG46） & 約870 & 約1.9 & 約0.13 & ― & ― & カ（熱-12） \\
人の体（平均） & ― & 約3.5 & ― & ― & ― & キ（熱-1） \\
\end{longtblr}\par}
{\footnotesize\color{bkGray}出典　ア：CRC Handbook の値（本文で照合していない\phyk）。イ：Incropera ほか『伝熱工学の基礎』の付表の代表値\phyk。ウ：数学シリーズ（6-3・8-4・8-13）と同じ数。エ：理科年表\phyk。オ：CoolProp 8.0（IAPWS-95）で計算。カ：規格・メーカー資料の代表値（JIS G 5501・JIS A 9521・作動油のメーカー資料）\phyk。キ：教科書・便覧の目安（人の体は 3.47 とする本が多い。土は 0.3〜0.8×10$^{-6}$ の中ほど）\phyk。\par
本文の水の比熱は 4.2（3けたの例題は 4.18。熱-1・熱-2・熱-6）。氷の融解熱は \qty{334}{kJ/kg}（エ。熱-1・熱-6・熱-10・熱-15）。鉄の熱拡散率は純鉄に近い値（数学 8-13）で、鋼の $k$ を使うと約 $1.4\times10^{-5}$。体膨張率 $\beta$（1/K、\qty{20}{\celsius}前後）：空気（圧力一定）$3.4\times10^{-3}$（$1/T$、熱-3）、ガソリン 約$9.5\times10^{-4}$（JIS K 2249 の体積換算表\phyk）、水 $2.1\times10^{-4}$（オ）。固体は $\beta\approx3\alpha$（熱-1 ⑦）。\par}

\phbkhead{9　熱の伝わりやすさの目安}
{\footnotesize
\begin{longtblr}{colspec={Q[l,m,wd=52mm] Q[l,m,wd=40mm] X[l,m]}, width=\linewidth, rowhead=1, \ph@tblstyle}
名前 & 値 & 出典・注意 \\
熱伝達率 $h$：自然対流の気体 & 2〜25\,W/(m$^2$·K) & Incropera ほか『伝熱工学の基礎』の表1.1（熱-12） \\
熱伝達率 $h$：強制対流の気体 & 25〜250\,W/(m$^2$·K) & 同じ \\
熱伝達率 $h$：自然対流の液体 & 50〜1000\,W/(m$^2$·K) & 同じ \\
熱伝達率 $h$：強制対流の液体 & 100〜\num{20000}\,W/(m$^2$·K) & 同じ \\
熱伝達率 $h$：沸騰・凝縮 & \num{2500}〜\num{100000}\,W/(m$^2$·K) & 同じ \\
1気圧の水の沸騰の極大熱流束 & 1.26〜2.09\,MW/m$^2$ & 抜山四郎、機械学会誌 37巻（1934年）。日本伝熱学会の公開資料で確かめた（熱-12） \\
壁の表面の熱抵抗 & 室内側 0.11、外気側 0.04\,m$^2$·K/W & 建築研究所「外皮の熱貫流率の計算」の資料（熱-11・熱-12） \\
複層ガラスの空気層（\qty{12}{mm}） & 0.17\,m$^2$·K/W & 伝導と放射からの見積もり 0.172（熱-3・熱-11・熱-12）。建築研究所の資料の 0.173 とも合う \\
窓の熱貫流率（計算） & 1枚ガラス 約6.5、複層ガラス 約3.1\,W/(m$^2$·K) & 熱-11 の計算値 \\
窓の熱貫流率（公表値） & 透明単板 5.8、透明複層 2.9\,W/(m$^2$·K) & 板硝子協会の表（熱-9・熱-11。表面の熱伝達の決め方がちがうので、計算値より少し小さい） \\
\end{longtblr}\par}

\phbkhead{10　放射率の目安（ふだんの温度）}
{\footnotesize
\begin{longtblr}{colspec={Q[l,m,wd=52mm] Q[c,m,wd=26mm] X[l,m]}, width=\linewidth, rowhead=1, \ph@tblstyle}
表面 & 放射率 $\varepsilon$ & 出典・注意 \\
人の皮膚・水・氷 & 0.95〜0.98 & Incropera の付録・日本機械学会『伝熱工学資料』の代表値の範囲\phyk（熱-13） \\
黒いペンキ・すす & 0.95前後 & 同じ \\
白いペンキ・白い紙 & 0.9前後 & 同じ（赤外線では「黒い」） \\
れんが・コンクリート・木・ガラス & 0.9前後 & 同じ \\
さびた鉄 & 0.6〜0.8 & 同じ \\
みがいたアルミ・アルミはく & 0.03〜0.05 & 同じ \\
銀めっき & 0.02前後 & 同じ（魔法びんの内側） \\
\end{longtblr}\par}
{\footnotesize\color{bkGray}放射率は、表面のよごれ・さび・あらさで大きく変わる。灰色体（波長によらない）として見た目安。\par}

\phbkhead{11　燃料の発熱量}
{\footnotesize
\begin{longtblr}{colspec={Q[l,m,wd=30mm] Q[l,m,wd=50mm] X[l,m]}, width=\linewidth, rowhead=1, \ph@tblstyle}
燃料 & 値 & 出典・注意 \\
メタン & 高発熱量 \qty{890.6}{kJ/mol}（約\qty{55.5}{MJ/kg}）、低発熱量 \qty{802.6}{kJ/mol} & \qty{25}{\celsius}。生成エンタルピーの差から（熱-1）。理科年表・JIS K 2279 との照合は\phyk \\
水素 & 高発熱量 \qty{285.8}{kJ/mol}（約\qty{142}{MJ/kg}）、低発熱量 \qty{241.8}{kJ/mol}（約\qty{120}{MJ/kg}） & NIST の生成エンタルピー（$-285.83$・$-241.83$\,kJ/mol）の差（熱-1・熱-10）\phyk \\
都市ガス 13A & 約\qty{45}{MJ}/m$^3_{\mathrm N}$（高発熱量） & ガス事業者の供給約款の標準熱量（熱-1・熱-0・熱-16） \\
炭素（黒鉛） & 燃焼熱 \qty{393.5}{kJ/mol}（約\qty{32.8}{MJ/kg}） & NIST の生成エンタルピー（CO$_2$ $-393.51$\,kJ/mol）（熱-16） \\
\end{longtblr}\par}
{\footnotesize\color{bkGray}高発熱量は、できた水蒸気が水にもどるときの潜熱（\qty{25}{\celsius}で \qty{44.0}{kJ/mol}）までふくめた値。m$^3_{\mathrm N}$ は\qty{0}{\celsius}・\qty{101.325}{kPa}での体積（熱-1）。発電所の効率は、日本では高発熱量を基準にすることが多い。\par}

\phbkhead{12　この巻で使う定数}
{\footnotesize
\begin{longtblr}{colspec={Q[l,m,wd=40mm] Q[l,m,wd=56mm] X[l,m]}, width=\linewidth, rowhead=1, \ph@tblstyle}
名前 & 値 & 出典・注意 \\
ボルツマン定数 $k_{\mathrm B}$ & \qty{1.380649e-23}{J/K}（決めた値） & 2019年の SI（BIPM の SI 文書 第9版） \\
アボガドロ定数 $N_{\mathrm A}$ & \qty{6.02214076e23}{mol^{-1}}（決めた値） & 同じ \\
モル気体定数 $R=N_{\mathrm A}k_{\mathrm B}$ & \qty{8.314462618}{J/(mol.K)} & 決めた値の積。本文では 8.31（数学 8-3 と同じ） \\
ステファン・ボルツマン定数 $\sigma$ & \qty{5.670374419e-8}{W/(m^2.K^4)} & 2019年の SI の決めた値（$k_{\mathrm B}$・プランク定数・光の速さ）から計算される値。CODATA 2018 と 2022 で同じ（熱-13・熱-14） \\
ウィーンの変位則の定数 & \qty{2.897771955e-3}{m.K} & 同じ \\
放射の第1・第2定数 & $c_1=\qty{3.741771852e-16}{W.m^2}$、$c_2=\qty{1.438776877e-2}{m.K}$ & 同じ \\
光の速さ $c_0$ & \qty{299792458}{m/s}（決めた値） & SI \\
電気素量（1\,eV） & \qty{1.602176634e-19}{J}（決めた値） & SI（熱-15） \\
ファラデー定数 & \qty{96485}{C/mol} & $N_{\mathrm A}$ と電気素量の積（熱-10） \\
\qty{0}{\celsius} & \qty{273.15}{K}（定義） & $\theta$/℃＝$T$/K$-273.15$（熱-1） \\
標準大気圧 & \qty{101325}{Pa}（決めた値） & 1気圧 \\
水の三重点 & \qty{273.16}{K}・\qty{611.657}{Pa} & IAPWS（熱-8・熱-10）。IF97 の原文の照合は\phyk \\
水の臨界点 & \qty{647.096}{K}（\qty{373.946}{\celsius}）・\qty{22.064}{MPa} & 同じ \\
二酸化炭素の三重点・臨界点 & \qty{216.59}{K}・\qty{0.518}{MPa}／\qty{304.13}{K}・\qty{7.377}{MPa} & Span・Wagner（1996）。CoolProp 8.0 の値（熱-10） \\
国際標準大気 & 海面 \qty{15}{\celsius}・\qty{101.325}{kPa}、\qty{11}{km}まで \qty{6.5}{K/km}で下がる（\qty{11}{km}で \qty{216.65}{K}） & ISO 2533（熱-3・熱-4） \\
乾いた空気が上がるときの温度の下がり方 & 約\qty{9.8}{K/km} & $g/c_p$ から（熱-3・熱-9） \\
太陽の有効温度・半径 & \qty{5772}{K}・\qty{6.957e8}{m} & IAU 2015 Resolution B3（熱-13） \\
太陽定数（地球の所の強さ） & \qty{1361}{W/m^2} & 同じ（太陽周期23の平均） \\
1天文単位 & \qty{1.495978707e11}{m}（決めた値） & IAU 2012 Resolution B2 \\
\end{longtblr}\par}

\phbkhead{13　実機の値と目安}
{\footnotesize
\begin{longtblr}{colspec={Q[l,m,wd=50mm] Q[l,m,wd=40mm] X[l,m]}, width=\linewidth, rowhead=1, \ph@tblstyle}
名前 & 値 & 出典・注意 \\
コンバインドサイクル発電の熱効率 & 63.08\,\%（西名古屋火力発電所7-1号、2018年） & GE のプレスリリース（熱-5・熱-7） \\
石炭火力（超々臨界圧）の熱効率 & 約40\,\%（送電端 約42\,\%、高位発熱量基準） & 経済産業省の技術評価資料（熱-8）。電力会社の公開値は\phyk（熱-5） \\
先進超々臨界圧（A-USC）の目標 & 蒸気\qty{700}{\celsius}級、送電端 46\,\% & 経済産業省・NEDO の資料（熱-8） \\
加圧水型原子炉の蒸気・熱効率 & 約\qty{6.5}{MPa}・約\qty{280}{\celsius}、約33\,\% & クルスコ原子力発電所の技術情報ほか（熱-5） \\
沸騰水型原子炉の主蒸気・復水器 & 約\qty{7}{MPa}・約\qty{286}{\celsius}、約\qty{5}{kPa} & 発電の教科書の典型値\phyk（熱-8） \\
乗用車のガソリンエンジンの最大の熱効率 & 40\,\%（ハイブリッド用 41\,\%） & トヨタの発表（2016年）（熱-7） \\
大型の船の2ストロークディーゼル & 50\,\% をこえる（低位発熱量基準） & 三菱重工の技術論文（熱-7） \\
軽油が自然に燃えだす温度 & 約\qty{250}{\celsius} & 数学 3-11 と同じ値。安全データシートの発火点は\phyk（熱-7） \\
家庭用エアコンの暖房の成績係数 & 約4〜6 & メーカーのカタログの値（熱-5）。JIS C 9612 の条件は\phyk \\
海洋温度差発電（久米島） & \qty{50}{kW}、表層 約\qty{26}{\celsius}・深層 約\qty{9}{\celsius} & 沖縄県の資料（熱-5） \\
人の発熱（座って軽い作業） & 顕熱 約\qty{70}{W}・潜熱 約\qty{45}{W} & ASHRAE Handbook の表\phyk（熱-9） \\
ガス給湯器の1号 & 約\qty{1.74}{kW}（水を\qty{25}{K}上げた湯を毎分\qty{1}{L}） & 日本ガス石油機器工業会の説明\phyk（熱-1） \\
旅客機の客室の気圧 & 約\qty{75}{kPa}より下げない（高さ約\qty{2400}{m}の大気と同じ程度） & \phyk（熱-3） \\
高圧蒸気滅菌（オートクレーブ） & \qty{121}{\celsius}で使うことが多い & 日本薬局方の参考情報\phyk（熱-10） \\
ダイヤモンドのモル比熱（\qty{298}{K}） & \qty{6.1}{J/(mol.K)} & NIST WebBook の 6.115（二次資料で確認）\phyk（熱-15） \\
窒素の標準モルエントロピー & \qty{191.6}{J/(mol.K)}（\qty{25}{\celsius}・\qty{1}{bar}） & 化学便覧・NIST の表\phyk（熱-15） \\
アルミ電解コンデンサーの寿命 & \qty{10}{K}上がると約半分 & 部品メーカーの技術資料の目安\phyk（熱-15） \\
スクラムジェットの実験機 & X-43A マッハ約9.6（2004年）、X-51A マッハ約5.1・約210秒（2013年） & NASA・アメリカ空軍の資料\phyk（熱-4） \\
発電所の温排水の温度差 & 約\qty{7}{K}（取水と放水の差の目安） & 電力会社の環境アセスメントの資料\phyk（熱-0・熱-8） \\
発電所の煙突から出るガス & 約\qty{100}{\celsius} & \phyk（熱-0） \\
太陽電池のモジュールの変換効率 & 約20\,\% & 市販品のふつうの値。メーカーの公開資料\phyk（熱-16） \\
設備利用率（1年の平均の出力÷定格） & 太陽光 15〜18\,\%（住宅の見本の計算は 15\,\%\phyk）、風力 約30\,\% & 発電コスト検証ワーキンググループ（2024年）の前提（熱-16） \\
発電コスト（2040年に新しくつくる発電所の試算） & 太陽光（事業用）7.0〜8.9 円/kWh、原子力 12.5 円/kWh 以上 & 同じ報告（資源エネルギー庁、2024年12月）。報道と解説資料で照合（熱-16） \\
日本の電気のうち火力の割合 & 約7割（2020年代のはじめ） & 資源エネルギー庁「総合エネルギー統計」\phyk（熱-16） \\
ITER（国際熱核融合実験炉） & 加熱 \qty{50}{MW}の10倍の \qty{500}{MW}が目標。重水素・三重水素の運転の開始 2039年 & ITER 機構の公開資料・記者会見（2024年7月）（熱-16） \\
\end{longtblr}\par}

\phbkhead{14　発電の方法と CO$_2$（ライフサイクル）}
{\footnotesize
\begin{longtblr}{colspec={Q[l,m,wd=46mm] Q[c,m,wd=34mm] X[l,m]}, width=\linewidth, rowhead=1, \ph@tblstyle}
発電の方法 & CO$_2$〔g/kWh〕（中央値） & 出典・注意 \\
石炭 & 820 & IPCC 第5次評価報告書 第3作業部会 Annex III の表 A.III.2（2014年）。原文の PDF には届かず、二次資料の表で照合（熱-16） \\
天然ガス（コンバインドサイクル） & 490 & 同じ \\
太陽光（大規模） & 48 & 同じ \\
地熱 & 38 & 同じ \\
水力 & 24 & 同じ \\
原子力 & 12 & 同じ \\
風力（陸上） & 11 & 同じ。陸上 11・洋上 12 とする表と、まとめて 11.5 とする表がある\phyk \\
\end{longtblr}\par}
{\footnotesize\color{bkGray}建てる・燃料をとる・こわすまでをふくめた、発電\qty{1}{kWh}あたりの温室効果ガス（CO$_2$ に換算した値）。\par}

\phbkhead{15　単位と換算}
{\footnotesize
\begin{longtblr}{colspec={Q[l,m,wd=30mm] Q[l,m,wd=40mm] X[l,m]}, width=\linewidth, rowhead=1, \ph@tblstyle}
量 & SI の単位 & ほかの単位との関係（単位つきの1で換算する。2-10） \\
温度 & K（ケルビン） & $\theta$/℃＝$T$/K$-273.15$。温度の差は\qty{1}{K}＝\qty{1}{\celsius}の差。華氏 $\theta_{\mathrm F}$/°F＝1.8\,$\theta$/℃＋32 \\
熱・エネルギー & J（ジュール） & 1\,cal＝\qty{4.184}{J}（熱化学カロリー。計量単位令の別表第六の十三。二次資料で確認、条文は\phyk。計量法では、食べ物の熱量と代謝の熱量にだけ使える）、1\,kWh＝\qty{3.6}{MJ}、1\,eV＝\qty{1.602e-19}{J} \\
熱の流れ・仕事率 & W（ワット）＝J/s & 1\,kcal/h≒\qty{1.163}{W} \\
エネルギーの統計の単位 & J & 1\,toe（石油換算トン）＝$10^7$\,kcal＝\qty{41.868}{GJ}（国際蒸気表カロリー 1\,kcal＝\qty{4.1868}{kJ}。IEA・OECD の定義）、1\,Btu≒\qty{1055}{J}（熱-16） \\
熱流束 & W/m$^2$ & 1\,kcal/(m$^2$·h)≒\qty{1.163}{W/m^2}（熱-12 の抜山の値） \\
圧力 & Pa（パスカル） & 1\,bar＝\qty{100}{kPa}、1気圧＝\qty{101.325}{kPa}、1\,Torr＝1\,mmHg≒\qty{133.3}{Pa}、1\,kgf/cm$^2$＝\qty{98.0665}{kPa} \\
比熱・比エントロピー & J/(kg·K) & 1\,kJ/(kg·K)＝1\,J/(g·K) \\
体積（燃料のガス） & m$^3$ & m$^3_{\mathrm N}$：\qty{0}{\celsius}・\qty{101.325}{kPa}にしたときの体積（熱-1） \\
\end{longtblr}\par}
{\footnotesize\color{bkGray}接頭語：G（ギガ、$10^9$）、M（メガ、$10^6$）、k（キロ、$10^3$）、m（ミリ、$10^{-3}$）、μ（マイクロ、$10^{-6}$）、n（ナノ、$10^{-9}$）。\par}

%% ---------------- なぜ？の道すじ ----------------
%% 修2-6（Q6-18）：35本が5ページめに1本だけはみ出していたので、この巻だけ、道すじの間（3.2→2.0 mm）と、問いと駅の間（1.2→0.8 mm）をつめる。
%%   形は共通部品（bk-physbody.sty の \phwhyroute）と同じ。共通部品は変えない
\renewcommand\phwhyroute[4]{%
  \par\addvspace{2.0mm}\noindent
  \begin{minipage}{\linewidth}
  \noindent{\ph@whyicon\hspace{0.5em}\sffamily\bfseries\color{bkInk}#1}\hfill{\sffamily\footnotesize\color{bkGray}入口：#2}\par\vspace{0.8mm}
  {\leftskip=1.6em\raggedright\noindent
   \foreach \s [count=\k] in {#3} {\ifnum\k>1 \ph@routeline\fi\ph@routestop{\s}}\par}
  \if\relax\detokenize{#4}\relax\else
  \vspace{0.4mm}{\leftskip=1.6em\noindent\sffamily\footnotesize\color{bkGray}使う数学：\foreach \m [count=\k] in {#4} {\ifnum\k>1 \hspace{0.25em}\fi\bk@pill{\m}}\par}\fi
  \end{minipage}\par}
\bkbacktitle{なぜ？の道すじ}
{\small 気になる「なぜ？」から、読む章の順番をたどる路線図です。札の横の数はこの巻のページです。ほかの巻の札（力・波・電・原）は、物理シリーズの『力学』『波動』『電磁気学』『原子』の章です。使う数学の章は、下の小さな札です。\par}
\vspace{1mm}
\phwhyroute{氷水は、氷がとけきるまで0\,℃のままなのはなぜ？}{熱-1}{熱-1}{2-3,2-16,6-3}
\phwhyroute{自転車の空気入れの筒の先は、なぜ熱くなる？}{熱-2}{力-4,熱-2,熱-3}{3-7,8-7}
\phwhyroute{高い山の上で、おかしの袋がふくらむのはなぜ？}{熱-3}{力-12,熱-3}{3-7,5-14}
\phwhyroute{ロケットのノズルは、なぜ一度細くしてから広げる？}{熱-4}{力-3,力-13,熱-3,熱-4}{5-6,8-8}
\phwhyroute{エンジンは、燃料の熱を全部、動きに変えられる？}{熱-5}{熱-2,熱-5}{3-7,8-7}
\phwhyroute{お湯と水をまぜたあと、なぜ元の温度に分かれない？}{熱-6}{熱-1,熱-5,熱-6,熱-15}{5-11,5-14}
\phwhyroute{ディーゼルエンジンは、なぜ火花なしで燃える？}{熱-7}{熱-3,熱-7}{3-11,8-7}
\phwhyroute{高い山の上では、なぜ水が100\,℃より低い温度でわく？}{熱-8}{熱-1,熱-8,熱-10}{3-5,3-12}
\phwhyroute{冷たいコップのまわりに、なぜ水がつく？}{熱-9}{熱-1,熱-8,熱-9,熱-11}{3-5,3-11,6-4}
\phwhyroute{圧力なべの中では、水は何℃でわく？}{熱-10}{熱-8,熱-10}{3-12,8-3}
\phwhyroute{二重窓にすると、なぜ部屋が冷えにくい？}{熱-11}{熱-11,熱-13}{6-4,8-13}
\phwhyroute{パソコンは、なぜファンで風を送って冷やす？}{熱-12}{力-14,熱-11,熱-12}{6-4,6-11}
\phwhyroute{たき火に向いた顔だけが熱くなるのはなぜ？}{熱-13}{熱-11,熱-13}{5-14,6-4}
\phwhyroute{空気の分子は、どれくらいの速さで飛んでいる？}{熱-14}{熱-3,熱-14}{3-17,8-6}
\phwhyroute{トランプは、切るほどなぜ順番がばらばらになる？}{熱-15}{熱-6,熱-15}{3-17,5-14}
\phwhyroute{発電所は、なぜ海の水をたくさん使う？}{熱-0}{熱-0,熱-5,熱-8}{3-7}
\phwhyroute{電気は、どの方法でつくるとよい？　何をくらべる？}{熱-16}{熱-0,熱-5,熱-6,熱-16}{5-12}
\phwhyroute{エアコンの暖房は、なぜ電気ストーブより電気代が安い？}{熱-5}{熱-2,熱-5,熱-9}{2-5}
\phwhyroute{冷蔵庫は、どうやって中の熱を外へ運び出す？}{熱-9}{熱-4,熱-5,熱-9}{2-5,3-5}
\phwhyroute{エアダスターを吹き続けると、なぜ缶が冷たくなる？}{熱-4}{熱-1,熱-4,熱-9}{2-3}
\phwhyroute{汗をかくと、なぜ体が冷える？}{熱-1}{熱-1,熱-9}{2-3}
\phwhyroute{熱気球は、なぜ空に浮かぶ？}{熱-3}{力-12,熱-3}{3-7}
\phwhyroute{雲の底は、どれくらいの高さにできる？}{熱-9}{熱-3,熱-9}{3-5}
\phwhyroute{晴れた冬の朝に、なぜ霜がおりる？}{熱-13}{熱-9,熱-13}{5-14}
\phwhyroute{魔法びんのお湯は、なぜ冷めにくい？}{熱-13}{熱-11,熱-12,熱-13}{6-4}
\phwhyroute{扇風機の風を当てると、お湯はどれだけ速く冷める？}{熱-12}{熱-1,熱-12}{3-12,6-4}
\phwhyroute{地面の下の温度は、なぜ季節でゆっくり変わる？}{熱-11}{熱-11}{6-4,9-1}
\phwhyroute{ドライアイスは、なぜ液体にならずに消える？}{熱-10}{熱-8,熱-10}{3-12}
\phwhyroute{氷の上は、なぜ滑りやすい？}{熱-1}{熱-1,熱-10}{2-3}
\phwhyroute{燃料電池の電圧は、何で決まる？}{熱-10}{熱-6,熱-10,電-4}{8-3}
\phwhyroute{太陽の表面の温度は、どうやってわかる？}{熱-13}{熱-13,原-2}{5-14}
\phwhyroute{地球の平均の気温は、何で決まる？}{熱-13}{熱-3,熱-13}{5-14}
\phwhyroute{ウランの濃縮では、どうやって軽い原子を集める？}{熱-14}{熱-14,原-15}{3-17}
\phwhyroute{電子部品は、なぜ熱に弱い？}{熱-15}{熱-12,熱-15}{5-14}
\phwhyroute{車はどうして動くの？}{熱-7}{力-2,熱-5,熱-7,電-12}{4-14,5-5,6-2}

%% ---------------- チェック表 ----------------
\bkbacktitle{できるようになったことチェック表}
{\small 各章の「ふりかえり」を集めた表です。章を学び終えたら1回め、次の巻に進む前に2回めのチェックをしよう。自信がない所は、ページの数字の所へもどろう。\par}\vspace{2mm}
% 自動生成（tools/mkchecklist.py p2）。各章の「ふりかえり」の項目を集めたもの（第13回）
\begingroup\footnotesize
\def\bkcb{\tikz[baseline=-0.5ex]\draw[\bkTF, line width=0.7pt, rounded corners=0.5pt] (0,-0.3em) rectangle (0.8em,0.5em);}
\begin{longtblr}{colspec={Q[c,m,wd=11mm] X[l,m] Q[c,m,wd=10mm] Q[c,m,wd=9mm] Q[c,m,wd=9mm]}, width=\linewidth, rowhead=1, colsep=3pt,
  row{1}={bg=bkVol!14, font=\sffamily\bfseries\scriptsize}, hlines={0.4pt,bkLine}, rowsep=0.7pt}
章 & できるようになったこと & ページ & 1回め & 2回め \\
{\sffamily\bfseries\color{\bkTT}熱-0} & 発電所に入る燃料の熱・出る電気・捨てる熱を、熱のバランス式と効率で数えられる。 & \pageref{ph:熱-0:sec2} & \bkcb & \bkcb \\
 & 何を無視したかを書いてから、捨てる熱を運ぶ海の水の量を、7つの段で求められる。 & \pageref{ph:熱-0:ex3} & \bkcb & \bkcb \\
 & 身近な「なぜ？」の答えが、この巻のどの章にあるかを探せる。 & \pageref{ph:熱-0:sec6} & \bkcb & \bkcb \\
{\sffamily\bfseries\color{\bkTT}熱-1} & 氷がとけている間、温度が変わらないわけを、潜熱で説明できる。 & \pageref{ph:熱-1:sec5} & \bkcb & \bkcb \\
 & $Q=mc\,\Delta T$ と熱のバランス式で、温まり方と混ぜたときの温度を求められる。 & \pageref{ph:熱-1:sec4} & \bkcb & \bkcb \\
 & 加熱曲線を読み、熱膨張の大きさを見積もれる。 & \pageref{ph:熱-1:sec6} & \bkcb & \bkcb \\
 & 燃料が燃えるときの空気の量と発熱量を、原子のバランスで求められる。 & \pageref{ph:熱-1:sec8} & \bkcb & \bkcb \\
{\sffamily\bfseries\color{\bkTT}熱-2} & 熱と仕事が同じエネルギーであることを、実験から説明できる。 & \pageref{ph:熱-2:sec1} & \bkcb & \bkcb \\
 & 第一法則 $\Delta U=Q+W$ で、エネルギーの出入りを数えられる。 & \pageref{ph:熱-2:sec4} & \bkcb & \bkcb \\
 & $p$-$V$ 図の面積から、気体がする仕事を求められる。 & \pageref{ph:熱-2:sec3} & \bkcb & \bkcb \\
 & 状態量と、道すじで変わる量を、全微分で見分けられる。 & \pageref{ph:熱-2:sec7} & \bkcb & \bkcb \\
{\sffamily\bfseries\color{\bkTT}熱-3} & ボイル・シャルルの法則と状態方程式で、気体の圧力・体積・温度を求められる。 & \pageref{ph:熱-3:sec2} & \bkcb & \bkcb \\
 & 定積・定圧・等温・断熱の4つの変化で、熱と仕事と内部エネルギーの出入りを数えられる。 & \pageref{ph:熱-3:sec4} & \bkcb & \bkcb \\
 & 断熱変化・ポリトロープ変化の式で、空気入れ・エンジン・圧縮機・大気の温度と圧力を計算できる。 & \pageref{ph:熱-3:sec6} & \bkcb & \bkcb \\
{\sffamily\bfseries\color{\bkTT}熱-4} & 定常流動のエネルギーの式で、タービン・圧縮機・ノズル・ドライヤーのエネルギーを計算できる。 & \pageref{ph:熱-4:sec3} & \bkcb & \bkcb \\
 & 音の速さ・マッハ数・よどみ点の温度を計算できる。 & \pageref{ph:熱-4:sec6} & \bkcb & \bkcb \\
 & 先細末広ノズルの、のどの状態と出口の面積を求められる。 & \pageref{ph:熱-4:sec8} & \bkcb & \bkcb \\
{\sffamily\bfseries\color{\bkTT}熱-5} & 熱機関の熱効率と、冷凍機・ヒートポンプの成績係数を、熱と仕事の矢印から計算できる。 & \pageref{ph:熱-5:sec2} & \bkcb & \bkcb \\
 & 第二法則の2つの言い方を説明し、カルノーの効率 $1-T_{\mathrm L}/T_{\mathrm H}$ を求められる。 & \pageref{ph:熱-5:sec3} & \bkcb & \bkcb \\
 & カルノーの定理から、熱機関と冷凍機の上限と、熱力学温度の考えを導ける。 & \pageref{ph:熱-5:sec5} & \bkcb & \bkcb \\
{\sffamily\bfseries\color{\bkTT}熱-6} & 熱÷温度 $\Delta S=Q/T$ を数えて、もどらない向きを言える。 & \pageref{ph:熱-6:sec2} & \bkcb & \bkcb \\
 & 水・氷・気体の $\Delta S$ を計算し、$T$-$s$ 線図の面積を読める。 & \pageref{ph:熱-6:sec7} & \bkcb & \bkcb \\
 & 生まれるエントロピーとエクセルギーで、機械の損失を数えられる。 & \pageref{ph:熱-6:sec9} & \bkcb & \bkcb \\
{\sffamily\bfseries\color{\bkTT}熱-7} & 4ストロークのエンジンの動きを $p$-$V$ 図で説明し、断熱圧縮のあとの温度を求められる。 & \pageref{ph:熱-7:sec1} & \bkcb & \bkcb \\
 & オットーサイクルとディーゼルサイクルの理論熱効率を、圧縮比から計算できる。 & \pageref{ph:熱-7:sec3} & \bkcb & \bkcb \\
 & ガスタービンの仕事比・再生・コンバインドサイクルの効率と、エンジンの出力を計算できる。 & \pageref{ph:熱-7:sec6} & \bkcb & \bkcb \\
{\sffamily\bfseries\color{\bkTT}熱-8} & 水のわく温度が、まわりの圧力（飽和蒸気圧）で決まることを説明できる。 & \pageref{ph:熱-8:sec1} & \bkcb & \bkcb \\
 & 乾き度と蒸気表（直線補間）で、湿り蒸気・過熱蒸気の比エンタルピーを求められる。 & \pageref{ph:熱-8:sec2} & \bkcb & \bkcb \\
 & ランキンサイクルの効率を計算し、再熱・再生の働きを説明できる。 & \pageref{ph:熱-8:sec6} & \bkcb & \bkcb \\
{\sffamily\bfseries\color{\bkTT}熱-9} & 湿度と露点を使って、コップや窓に水がつくわけを説明できる。 & \pageref{ph:熱-9:sec6} & \bkcb & \bkcb \\
 & 絶対湿度・比エンタルピー・空気線図で、エアコンが空気から取る水と熱を計算できる。 & \pageref{ph:熱-9:sec5} & \bkcb & \bkcb \\
 & 冷凍サイクルの成績係数（COP）を、冷媒の表から計算できる。 & \pageref{ph:熱-9:sec2} & \bkcb & \bkcb \\
{\sffamily\bfseries\color{\bkTT}熱-10} & 蒸気表とファンデルワールスの式で、わく温度と実在気体の圧力を求められる。 & \pageref{ph:熱-10:sec8} & \bkcb & \bkcb \\
 & $\ln p$ と $1/T$ のグラフから蒸発熱を読み、わく温度を見積もれる。 & \pageref{ph:熱-10:sec3} & \bkcb & \bkcb \\
 & 自由エネルギーとマクスウェルの関係式から、相平衡の式を導ける。 & \pageref{ph:熱-10:sec7} & \bkcb & \bkcb \\
{\sffamily\bfseries\color{\bkTT}熱-11} & 熱の伝わり方（伝導・対流・放射）を、身近な例で見分けられる。 & \pageref{ph:熱-11:sec1} & \bkcb & \bkcb \\
 & フーリエの法則で、平らな壁を通りぬける熱の流れを求められる。 & \pageref{ph:熱-11:sec2} & \bkcb & \bkcb \\
 & 熱抵抗の直列と熱貫流率で、窓や壁から逃げる熱と、表面の温度を求められる。 & \pageref{ph:熱-11:sec4} & \bkcb & \bkcb \\
 & 壁を小さな箱に分けた表計算で、壁の中の温度の変わり方を追える。 & \pageref{ph:熱-11:sec5} & \bkcb & \bkcb \\
 & 管の保温材の熱抵抗と、ビオ数・時定数を使って、冷え方を見積もれる。 & \pageref{ph:熱-11:sec8} & \bkcb & \bkcb \\
 & 熱方程式を導き、板の冷え方や地中の温度を、変数分離の答えで読める。 & \pageref{ph:熱-11:sec9} & \bkcb & \bkcb \\
{\sffamily\bfseries\color{\bkTT}熱-12} & ニュートンの冷却則と熱伝達率の目安で、物から流れへ逃げる熱を見積もれる。 & \pageref{ph:熱-12:sec2} & \bkcb & \bkcb \\
 & 冷めていく温度の測定から、表計算で熱伝達率を求められる。 & \pageref{ph:熱-12:sec3} & \bkcb & \bkcb \\
 & フィンを小さな箱に分けた表計算で、フィンの温度と効きめを求められる。 & \pageref{ph:熱-12:sec4} & \bkcb & \bkcb \\
 & 無次元数と相関式を選んで、強制対流・自然対流の熱伝達率を求められる。 & \pageref{ph:熱-12:sec6} & \bkcb & \bkcb \\
 & 沸騰曲線と限界熱流束、凝縮の熱伝達を説明し、大きさを見積もれる。 & \pageref{ph:熱-12:sec9} & \bkcb & \bkcb \\
{\sffamily\bfseries\color{\bkTT}熱-13} & 物が出す放射の強さ $\sigma T^4$ と、いちばん強い波長を計算できる。 & \pageref{ph:熱-13:sec2} & \bkcb & \bkcb \\
 & 放射率と形態係数を使って、2つの面の間で放射で移る熱を計算できる。 & \pageref{ph:熱-13:sec6} & \bkcb & \bkcb \\
 & 対数平均温度差と温度効率で、熱交換器の熱の量と大きさを求められる。 & \pageref{ph:熱-13:sec8} & \bkcb & \bkcb \\
{\sffamily\bfseries\color{\bkTT}熱-14} & 分子が壁にぶつかる力積から、気体の圧力と温度の意味を説明できる。 & \pageref{ph:熱-14:sec2} & \bkcb & \bkcb \\
 & 気体の内部エネルギーと比熱を、分子の運動から計算できる。 & \pageref{ph:熱-14:sec4} & \bkcb & \bkcb \\
 & 速さの分布と平均自由行程から、粘度・熱伝導率・拡散を見積もれる。 & \pageref{ph:熱-14:sec5} & \bkcb & \bkcb \\
{\sffamily\bfseries\color{\bkTT}熱-15} & 並び方の数を数えて、変化がなぜ一方の向きにだけ進むのかを説明できる。 & \pageref{ph:熱-15:sec2} & \bkcb & \bkcb \\
 & $S=k_{\mathrm B}\ln\Omega$ で、並び方の数からエントロピーを計算できる。 & \pageref{ph:熱-15:sec4} & \bkcb & \bkcb \\
 & ボルツマン分布と分配関数から、2準位の系や固体の比熱を導ける。 & \pageref{ph:熱-15:sec7} & \bkcb & \bkcb \\
{\sffamily\bfseries\color{\bkTT}熱-16} & \qty{1}{kWh}の電気に要る燃料の熱と、燃やして出る二酸化炭素の量を、効率と原子の数のバランスから見積もれる。 & \pageref{ph:熱-16:sec2} & \bkcb & \bkcb \\
 & 発電や給湯の方法を、効率・温度・設備利用率・一生で出る二酸化炭素など、いくつもの物差しでくらべられる。 & \pageref{ph:熱-16:sec9} & \bkcb & \bkcb \\
 & ヒートポンプや燃料電池が得をするわけを、エクセルギー（第二法則）で説明できる。 & \pageref{ph:熱-16:sec6} & \bkcb & \bkcb \\
\end{longtblr}
\endgroup


%% ---------------- つぎの巻へ ----------------
\bkbacktitle{つぎの巻へ}
\newcommand\bknextpill[1]{\bk@pill{#1}}
{\small 熱学で、熱と仕事を数え、熱機関・冷凍機・伝熱・分子の運動を調べる道具がそろいました。物理シリーズの『波動』では、ゆれと波を調べます。熱学で導いた音速や、ふく射の光、部屋の温度の制御が、そこで生きてきます。\par}\vspace{3mm}
\begin{tcolorbox}[enhanced, colframe=phVol3, colback=white, boxrule=0.9pt, arc=2.5mm, left=3mm, right=3mm, top=2mm, bottom=2mm,
  title={物理シリーズの『波動』：ゆれと波を読む}, fonttitle=\sffamily\bfseries, colbacktitle=phVol3, coltitle=white]
\small
\begin{itemize}[leftmargin=1.2\zw, itemsep=2pt, topsep=1pt, label={\color{phVol3}●}]
\item \bk@pill{波-4}\ ねらいどおりに動かす：フィードバック制御（熱-9 のエアコンで、熱-11 の熱容量・熱抵抗の部屋を一定の温度に保つ）
\item \bk@pill{波-9}\ 音の大きさをはかる：音速・デシベル・騒音（熱-4 で導いた気体の音速 $a=\sqrt{\gamma R_{\mathrm s}T}$ を、液体・固体へ広げる）
\item \bk@pill{波-13}\ 光の波：干渉・回折・偏光・散乱（熱-13 のふく射の波長と色）
\end{itemize}
\end{tcolorbox}
\phbkhead{この巻の道具は、この先でこう使う}
{\small この巻の章と、その道具を使う物理シリーズのほかの巻（『波動』から先）の章です（章のはじめの「道しるべ」の「この章の先に」から自動で集めた表）。\par}\vspace{2mm}
% 自動生成（tools/mknext.py p2）。各章の道しるべ「この章の先に」から、後ろの巻の章を集めたもの（第13回）
\begin{longtblr}{colspec={Q[l,m,wd=58mm] X[l,m]}, width=\linewidth, rowhead=1, row{1}={bg=bkVol!14, font=\sffamily\bfseries\footnotesize},
  hlines={0.4pt,bkLine}, rowsep=1.5pt, cells={font=\footnotesize}}
この巻の章 & この道具を使う章 \\
{\bknextpill{熱-3}\ ふくらむ・縮む気体：理想気体の状態変化} & {\bknextpill{波-9}\ 音の大きさをはかる：音速・デシベル・騒音} \\
{\bknextpill{熱-4}\ 流れる気体の熱：エンタルピーとノズル} & {\bknextpill{波-9}\ 音の大きさをはかる：音速・デシベル・騒音} \\
{\bknextpill{熱-7}\ エンジンの中のサイクル：ガスサイクル} & {\bknextpill{電-12}\ モーターと変圧器のしくみ：電気機器} \\
{\bknextpill{熱-9}\ 冷やす・湿り気をはかる：冷凍と湿り空気} & {\bknextpill{波-4}\ ねらいどおりに動かす：フィードバック制御} \\
{\bknextpill{熱-10}\ わく温度を圧力で読む：自由エネルギーと相平衡} & {\bknextpill{電-4}\ 電気の流れと熱：電流と抵抗} \\
{\bknextpill{熱-11}\ 壁を通りぬける熱：熱伝導と熱抵抗} & {\bknextpill{波-4}\ ねらいどおりに動かす：フィードバック制御} \\
{\bknextpill{熱-13}\ 光で伝わる熱：ふく射と熱交換器} & {\bknextpill{波-12}\ 光の進み方：反射・屈折・レンズ\\ \bknextpill{波-13}\ 光の波：干渉・回折・偏光・散乱\\ \bknextpill{原-2}\ 光はつぶでもある：光電効果と光子} \\
{\bknextpill{熱-14}\ 温度の正体：気体分子運動論} & {\bknextpill{波-9}\ 音の大きさをはかる：音速・デシベル・騒音\\ \bknextpill{原-6}\ 壁をすりぬける：トンネル効果と調和振動子} \\
{\bknextpill{熱-15}\ 並び方を数える：統計力学の入口} & {\bknextpill{原-4}\ 原子が出す光：エネルギー準位とスペクトル\\ \bknextpill{原-6}\ 壁をすりぬける：トンネル効果と調和振動子\\ \bknextpill{原-8}\ 固体の中の電子：エネルギーバンド\\ \bknextpill{原-9}\ 一方通行の電気：pn接合とLED・太陽電池} \\
{\bknextpill{熱-16}\ エネルギーをどう選ぶか} & {\bknextpill{電-16}\ 電気を届ける・電気を変える\\ \bknextpill{原-15}\ 原子炉の中の中性子：核分裂と原子炉\\ \bknextpill{原-16}\ 原子の力とつきあう} \\
\end{longtblr}


%% ---------------- さくいん ----------------
\bkbacktitle{さくいん}
{\small 数字はページです。その語がはじめて出てくるページや、詳しく説明しているページをのせています。}\par\vspace{2mm}
\raggedcolumns\begin{multicols}{2}\small\setlength{\parskip}{0pt}
\directlua{bk.printidx()}
\end{multicols}

%% ---------------- 奥付 ----------------
%% 第21回（R4-43）：さくいんのあとのメモのページの柱を「メモ」にする
\clearpage\bksetrunning{}{メモ}{メモ}%
\bkpadfour
\global\bk@decofalse\bkfolio{false}
\null\vfill
\noindent\begin{minipage}{\linewidth}\small
{\sffamily\bfseries\Large\color{bkDeep}\phSeries　第2巻}\par\vspace{1mm}
{\sffamily\bfseries\large 熱学：エネルギーを使いこなす}\par\vspace{3mm}
{\color{bkVol}\hrule height 1pt}\vspace{3mm}
\begin{tabular}{@{}l@{\hspace{6mm}}l@{}}
編者 & \phEditor \\
 & 委員長　小学校長（まとめ役） \\
 & 委員　数学・電気工学・機械工学・教育学・量子力学・建築工学・\\
 & \phantom{委員　}半導体工学の大学教授、幼稚園教諭、塾講師、\\
 & \phantom{委員　}特別支援学校教諭、心理士、科学解説者、小学校教諭 \\
デザイン & ユニバーサルデザイン担当委員 \\
挿絵 & イラスト担当委員 \\
先に読む本 & 「物理を理解するための数学」全9巻、物理シリーズ第1巻『力学』 \\
版 & 第1版 Ver1（第23回編集委員会）　2026年11月 \\
\end{tabular}\par\vspace{3mm}
{\color{bkVol}\hrule height 1pt}\vspace{2mm}
{\footnotesize この本の問題の答えは、巻末の「こたえ」にあります。詳しい解き方は、別冊の問題集と指導書にのせる予定です。蒸気表・空気線図は、IAPWS-IF97 の式と CoolProp 8.0 から計算しました。}
\end{minipage}
\makeatother

